\chapter{Versuchsaufbau}
\label{ch:versuchsaufbau}

Untersucht wird, ob und wie stark die Netzwerktopologie beeinflusst und wie
schnell und wie gut ein Verteidigungsagent lernt. Die Topologie ist die
unabhängige Variable, jedoch nicht der einzige variierte Faktor: Der
Verteidiger trainiert stets gegen einen bestimmten eingefrorenen Angreifer.
Ob und in welchem Umfang dessen Herkunft die Aufgabe beeinflusst, ist offen und
selbst Teil der Untersuchung. Beide Größen werden deshalb systematisch
gekreuzt, sodass sich ihre Beiträge trennen lassen (\Cref{sec:phasen}).

Alles Übrige muss konstant gehalten werden, sonst lassen sich beobachtete
Unterschiede nicht mehr zuordnen. Bevor dieser Aufbau im Einzelnen beschrieben
wird, ist jedoch ein Umweg nötig: Der hier verwendete Reward und Aktionsraum
sind nicht die von MARLon übernommenen, sondern eine eigene Überarbeitung.
\Cref{sec:vorarbeiten} begründet, warum diese Überarbeitung nötig wurde.

\section{Von MARLon zum eigenen Aufbau}
\label{sec:vorarbeiten}

Dieser Abschnitt beschreibt, wie aus dem übernommenen Rahmenwerk der in
dieser Arbeit verwendete Aufbau wurde. Er benennt zunächst den geerbten
Ausgangspunkt, schildert dann einen vollständigen Durchgang durch die
Versuchsmatrix und die drei Fehler, die dieser sichtbar machte, und begründet
daraus die anschließende Überarbeitung.

\subsection{Der geerbte Ausgangspunkt}
\label{sec:marlon-original}

MARLon~\autocite{marlon} stellt die Trainingsumgebung für beide Agenten
bereit, wie in \Cref{sec:cbs} beschrieben. Für den Verteidiger legt die
Originalarbeit dabei eine bestimmte Reward-Struktur fest: Sein Reward ergibt
sich durch Negation des jeweils letzten Angreifer-Rewards, ergänzt um eine
große Strafe bei Unterschreiten der Netzverfügbarkeit. Die Autoren begründen
diese Wahl damit, dass der Verteidiger den Fortschritt des Angreifers
minimieren solle, und halten als unmittelbare Folge fest, dass sein maximal
erreichbarer Reward damit bei null liegt und mit jedem Fortschritt des
Angreifers weiter sinkt~\autocite{marlon}. Prävention ist in dieser Formel nicht vorgesehen, und einen Anreiz zur
Rückeroberung enthält sie ebenso wenig: \ac{CBS} schreibt den Knotenwert nur bei der ersten Übernahme gut, sodass
ein späterer Besitzwechsel desselben Knotens den Reward beider Seiten
unverändert lässt. Gefördert wird damit gerade nicht das Verhalten, auf das
es in dieser Arbeit ankommt: den Angreifer am Vordringen zu hindern und
bereits verlorene Systeme zurückzugewinnen.

Für den Aktionsraum des Verteidigers sieht MARLon fünf Aktionsarten vor:
%
\begin{enumerate}
  \item einen Knoten neu aufsetzen,
  \item Verkehr sperren,
  \item Verkehr freigeben,
  \item einen Dienst stoppen,
  \item einen Dienst wieder starten.
\end{enumerate}
%
Port und Richtung sind dabei zusätzliche Parameter, die der Agent selbst
wählt. Für den Angreifer
identifiziert die Originalarbeit den Umgang mit ungültigen Aktionen als eigenes
Problem. Gültig ist eine Aktion, wenn der Angreifer sie im aktuellen Zustand
überhaupt versuchen kann: Der Ausgangsknoten muss übernommen, das Ziel entdeckt
und die eingesetzten Zugangsdaten müssen bereits erbeutet sein. Ob der Versuch
gelingt, entscheidet sich unabhängig davon. Ein Verbindungsversuch auf einen
Port, auf dem das Ziel keinen Dienst anbietet, ist eine gültige Aktion, die
lediglich folgenlos bleibt. Eine Bestrafung ungültiger Aktionen würde dazu
führen, dass der Agent riskante, aber
notwendige Aktionen ganz meidet; die Autoren entscheiden sich stattdessen für
Reward~0 bei ungültigen Aktionen, kombiniert mit einer zufälligen
Ersatzwahl. Diese Entscheidung ist im Code der Umgebung selbst verankert, nicht
nur im Trainingsablauf.

\subsection{Eigene Vorarbeiten und der Referenzlauf}
\label{sec:referenzlauf}

Dass der Verteidiger damit weder für verhinderte noch für rückgängig
gemachte Übernahmen etwas erhält, war der Anlass, seinen Reward in mehreren
Schritten zu überarbeiten. Ein zwischenzeitlicher
Entwurf war ereignisbasiert und vergütete jedes Neuaufsetzen eines
kompromittierten Knotens mit einem festen Betrag. Damit belohnte er einen
beliebig wiederholbaren Vorgang, und
zwar einseitig: \ac{CBS} schreibt einen Knotenwert nur bei der ersten Übernahme
gut, sodass dem Angreifer eine Rückeroberung nichts einbrachte, während der
Verteidiger für jedes Neuaufsetzen desselben Knotens erneut vergütet wurde. Ein
einzelner Knoten wechselte dadurch fortwährend zwischen kompromittiert und neu
aufgesetzt, ohne dass sich die Lage im Netz über einen solchen Zyklus hinweg
verändert hätte, und der Verteidiger wurde für jede Runde erneut belohnt. Das ist genau
der in \Cref{sec:rl-schwierigkeiten} beschriebene Fall, in dem ein Agent formal
korrekt maximiert und inhaltlich nichts erreicht. Daraufhin wurde zur Kopplung
aus \Cref{sec:marlon-original} zurückgekehrt, die MARLon ursprünglich vorsieht.

Über diese Kopplung hinaus umfasste der Aufbau zu diesem Zeitpunkt vier
eigene Ergänzungen:
%
\begin{enumerate}
  \item vier Vergleichstopologien anstelle des einen von MARLon
        vorgesehenen festen Netzes (\Cref{sec:topomodell}),
  \item ein Invalid-Action-Masking für den Angreifer nach Huang und
        Ontañón~\autocite{huang2022masking}, das MARLons
        Zufallsersatzverfahren ablöst (\Cref{sec:masking}),
  \item ein Haltebonus je gehaltenem Knoten und Zeitschritt
        (\Cref{sec:rewards}),
  \item ein Abzug in Höhe des Knotenwerts, wenn ein gehaltener Knoten neu
        aufgesetzt wird (\Cref{sec:rewards}).
\end{enumerate}
%
Die zweite Ergänzung verdient eine Erläuterung: Statt eine ungültige Aktion
zu raten und auf eine zufällige gültige auszuweichen, wählt der Angreifer
ausschließlich aus der im jeweiligen Zustand gültigen Menge. Wirkungslose
Aktionen entfielen dadurch praktisch vollständig, und der Angreifer
erreichte das Crown Jewel zuverlässig, statt wie zuvor nur in
Ausnahmefällen. Erst das machte ihn als Gegner brauchbar, denn die
Lerneffizienz eines Verteidigers lässt sich nur gegen einen Angreifer
messen, der überhaupt etwas ausrichtet.

Von der übernommenen Reward-Struktur blieben damit zwei Bestandteile unberührt,
und gerade sie bestimmen die erreichbare Obergrenze: die Kopplung des
Verteidiger-Rewards an den negierten Angreifer-Reward und die Begrenzung des
Angreifer-Schrittwerts nach unten auf null. Letztere ist in \ac{CBS} selbst
verankert und wurde bewusst beibehalten; \Cref{sec:rewards} legt dar, weshalb.
Die eigenen Ergänzungen verschieben den Angreifer-Reward, heben diese
Obergrenze aber nicht auf.

Auf diesem Stand entstand der im Folgenden als \emph{Referenzlauf} bezeichnete
Durchgang durch die vollständige Versuchsmatrix. Er ist der Vergleichsmaßstab
dieser Arbeit: Sein Ergebnis begründet die anschließend vorgenommenen
Änderungen am Aufbau, und an ihm wird deren Wirkung gemessen. Die Schwächen,
die die folgenden Abschnitte beschreiben, wurden erst durch ihn sichtbar und
waren zum Zeitpunkt seiner Durchführung nicht bekannt. Darin unterscheidet er
sich von den übrigen Zwischenständen des Aufbaus, die nur dort zur Sprache
kommen, wo sie eine Entwurfsentscheidung erklären. Die Matrix kreuzt die vier
Verteidigungstopologien mit sechs Angreifern: fünf davon wurden auf jeweils
genau einer Topologie trainiert, ein sechster nacheinander auf allen. Wegen der
in \Cref{sec:rl-schwierigkeiten} genannten Stochastik wird jede der
24~Kombinationen fünfmal wiederholt, jeweils mit einem anderen Startwert für die
Zufallsgeneratoren, einem sogenannten \emph{Seed}. Das ergibt
\num{120}~Trainingsläufe.

Als Budget waren dabei je \num{500000} Schritte vorgesehen, ausgeschöpft wurde
es aber nicht durchgängig: Das Konvergenzkriterium aus \Cref{sec:konvergenz}
diente im Referenzlauf noch als Abbruchbedingung und beendete das Training,
sobald es erfüllt war. Für die später berichteten Läufe wurde davon abgerückt;
\Cref{sec:konvergenz} begründet das. Die vier Ergänzungen sind
auch im überarbeiteten Aufbau enthalten und werden in den folgenden
Abschnitten beschrieben; was erst nach dem Referenzlauf hinzukam, ist
Gegenstand von \Cref{sec:konsequenz}.

Das Ergebnis bestätigte, was \Cref{sec:marlon-original} bereits aus der Formel
ableitet: In allen \num{8642} protokollierten Trainingsepisoden war der
Verteidiger-Reward negativ, der höchste beobachtete Einzelwert betrug
$-14{,}0$. Das ist keine Fehlfunktion, sondern die Obergrenze null, empirisch
bestätigt. Für eine Metrik, die ursprünglich vorgesehen war, hat das
unmittelbare Folgen: Break-even, definiert als die Episode, ab der der
geglättete Reward nicht mehr negativ ist, kann unter dieser Kopplung
strukturell nie eintreten und entfällt deshalb vollständig
(\Cref{sec:metriken}).

\subsection{Drei Fehler in der Umsetzung}
\label{sec:marlon-fehler}

Über die reine Reward-Formel hinaus fanden sich drei Stellen, an denen die
Sperre als Mittel der Prävention beeinträchtigt war. Sämtliche drei wurden erst
nach dem Referenzlauf behoben und waren während seiner gesamten Dauer wirksam.

Der erste betrifft die Liste der sperrbaren Ports und ist der eigenen Arbeit
zuzurechnen. MARLon hält sie fest verdrahtet im Code, passend zu der
Umgebung, für die das Rahmenwerk entworfen wurde; dort ist sie zutreffend. Mit
den in \Cref{sec:topomodell} beschriebenen Topologien kamen jedoch andere
Dienste ins Netz, und die Liste wurde dabei nicht mitgeführt. Von den sechs
tatsächlich vorkommenden Ports fehlten ihr zwei, während sie zwei weitere
enthielt, die in keiner der Topologien auftreten. Der Datenbank-, der Backup-
und der Dateiserver, darunter das Crown Jewel, ließen sich dadurch nie
ansprechen. Für die späteren Läufe wird die Liste nicht mehr fest vorgegeben,
sondern aus der Umgebung abgeleitet und folgt damit jeder Topologie.

Die beiden folgenden Punkte stammen dagegen aus dem unveränderten
MARLon-Code. Sie blieben im Referenzlauf ohne Folgen und wurden erst durch
die spätere Überarbeitung zum Hindernis; der Vollständigkeit halber gehören sie
dennoch hierher.

Die Methode zum Sperren von Verkehr entfernte lediglich eine vorhandene
Freigabe-Regel, statt eine ausdrückliche Sperr-Regel zu setzen. Unterbunden war
der Verkehr danach tatsächlich, weil \ac{CBS} einen Port ohne Regel als gesperrt
behandelt. Ein so gesperrter Port sah anschließend jedoch genauso aus wie einer,
für den nie eine Freigabe bestand, und das ist in den hier verwendeten
Topologien der bauliche Normalzustand: Sie arbeiten mit Freigabelisten, sodass
die meisten Ports ohnehin keine Regel tragen. Ein Eingriff des Verteidigers war
damit nicht mehr von der Ausgangslage zu unterscheiden. Für den Angreifer
änderte das nichts, denn er stieß in beiden Fällen auf eine geschlossene Tür.
Sobald jedoch der Reward eine bewusst gesetzte Sperre erkennen und belohnen
soll, wie in \Cref{sec:rewards} beschrieben, ist die Unterscheidung
unentbehrlich. Die zugehörige Methode zum Freigeben von Verkehr war zudem
fehlerhaft: Beide Zweige der Fallunterscheidung schrieben in die eingehende
Regelliste, sodass sich eine ausgehende Sperre nicht wieder aufheben ließ.

Schließlich zählte die Verfügbarkeitsberechnung nur, ob ein Dienst läuft, nicht
aber, ob sein Port gesperrt ist. Sperren kostete damit keine Verfügbarkeit. Im
Referenzlauf fiel auch das nicht ins Gewicht, weil eine Sperre dort ohnehin
nichts einbrachte. Wird sie dagegen belohnt, entsteht daraus ein Fehlanreiz:
Der Verteidiger könnte das Netz vollständig abriegeln, ohne dafür je das
\ac{SLA} zu verletzen.

Für den Referenzlauf bleibt damit vor allem der erste Punkt. Er allein erklärt
noch nicht, warum der Verteidiger ausschließlich das Neuaufsetzen lernte, denn
die entscheidende Ursache liegt in der Reward-Formel: Ein verhinderter Zugriff
bringt dem Angreifer null ein und dem Verteidiger über die Kopplung ebenfalls
null. Die fehlerhafte Portliste kam hinzu und schloss die wertvollsten Systeme
von vornherein aus. Und die einzige verbliebene präventive Handlung wurde sogar
bestraft: Das Stoppen eines Dienstes auf einem nicht kompromittierten Knoten
kostete zehn Punkte, und da sich nicht vorhersehen lässt, welchen Knoten der
Angreifer als Nächstes angreift, traf diese Strafe den Verteidiger regelmäßig.
Von den präventiven Mitteln blieb ihm keines, das sich gelohnt hätte.

\subsection{Konsequenz für diese Arbeit}
\label{sec:konsequenz}

Aus diesem Befund folgt nicht, dass sich mit der Fragestellung nach dem
Einfluss der Topologie nichts anfangen ließe. Der Referenzlauf bleibt
verwertbar, allerdings für eine engere Aussage, als ursprünglich
beabsichtigt: Er beschreibt einen Verteidiger, der ausschließlich reaktiv
handeln kann, also Wiederherstellung statt Prävention. Reimaging ist selbst
eine Form der Abwehr, denn es entzieht dem Angreifer den Zugang und macht
dessen Gewinn rückgängig, auch wenn es keinen Angriff verhindert.

Für die in dieser Arbeit berichteten Ergebnisse wurde der Aufbau aus zwei Gründen
grundlegend überarbeitet: Erstens lassen sich die drei in
\Cref{sec:marlon-fehler} genannten Punkte als Implementierungsfehler und nicht
als Entwurfsentscheidungen einordnen; sie zu beheben ist unabhängig von der
Forschungsfrage geboten. Zweitens lässt erst ein Verteidiger, der Prävention
tatsächlich ausführen und dafür belohnt werden kann, die Frage nach dem
Einfluss der Topologie auf einen \emph{präventiven} Verteidiger überhaupt
zu; unter dem Referenzlauf wäre jede Topologie mit derselben Unmöglichkeit
konfrontiert gewesen. Die folgenden Abschnitte beschreiben den überarbeiteten
Aufbau, wie er den in \Cref{ch:ergebnisse} berichteten Läufen zugrunde liegt.
Wo eine Änderung das Ergebnis unmittelbar beeinflusst, wird auf den Referenzlauf
als Vergleichspunkt zurückverwiesen.

\section{Wahl der Simulationsumgebung und des Lernverfahrens}
\label{sec:umgebungswahl}

Zwei Entscheidungen liegen dem Versuch zugrunde, bevor sich Fragen seiner
Ausgestaltung stellen: in welcher Simulationsumgebung er stattfindet und
mit welchem Lernverfahren die Agenten trainiert werden. Dieser Abschnitt
begründet beide.

\subsection{Die Simulationsumgebung}

Von den in \Cref{sec:umgebungen} vorgestellten Umgebungen fiel die Wahl auf \ac{CBS} in Verbindung mit MARLon. Maßgeblich dafür sind drei
Anforderungen, die sich unmittelbar aus der Fragestellung ergeben: eine frei
definierbare Netzwerktopologie, ein lernender Agent auf beiden Seiten und ein
Durchsatz, der eine Versuchsmatrix dieses Umfangs erlaubt. Die beiden
Alternativen scheitern jeweils an mehreren dieser Anforderungen.

Zu beachten ist dabei, dass alle drei Umgebungen fortlaufend weiterentwickelt
werden. Der folgende Vergleich bezieht sich auf den Stand, den die jeweiligen
Veröffentlichungen dokumentieren; spätere Erweiterungen sind, soweit sie belegt
sind, angeführt.

Die erste Anforderung ist die entscheidende, weil sie den Untersuchungsgegenstand
selbst betrifft. Dass \ac{CBS} das Netz unmittelbar im Programmcode als
Graph definiert (\Cref{sec:umgebungen}), ermöglicht erst den in
\Cref{sec:umgebung}
beschriebenen Aufbau, in dem sich die vier Topologien allein in ihrer
Segmentierung unterscheiden, während Knotenzahl und Knotenwerte konstant
bleiben. Die dokumentierten Szenarien der beiden anderen Umgebungen sind
demgegenüber klein und auf andere Zwecke zugeschnitten. Eine Topologie zu
verändern bedeutet dort, die zugrunde liegende Infrastruktur umzubauen.

Die zweite Anforderung betrifft die Gegenseite. Der Verteidiger soll sich gegen
ein Angriffsverhalten behaupten, das seinerseits aus Erfahrung entstanden ist
und nicht aus einer festen Regel. Wie \Cref{sec:beidseitig} zeigt, ist der
lernende Verteidiger in beiden Alternativen jedoch nicht umgesetzt. Genau diese
Lücke schließt MARLon, indem es für beide Rollen \ac{RL}-Agenten bereitstellt.

Die dritte Anforderung ist eine praktische. Die in dieser Arbeit ausgewerteten
120 Trainingsläufe umfassen zusammen rund 60~Millionen
Simulationsschritte.
Rechnet man die für CyGIL dokumentierten Episodendauern
(\Cref{sec:umgebungen}) auf diesen Umfang hoch, liegt das Ergebnis jenseits
des für eine Abschlussarbeit verfügbaren Zeitrahmens; der dort ebenfalls
genannte Simulator CyGIL-S nimmt diese Hürde nicht, weil er den Betrieb des
Emulators voraussetzt. CybORG bietet zwar
einen reinen Simulationsmodus, betreibt seinen Emulator jedoch über eine
Cloud-Infrastruktur~\cite{standen2021cyborg}, was neben dem Zeit- auch einen
Kostenaufwand nach sich zieht.

Mit der Wahl wird zugleich eine Einschränkung übernommen: Als reine Simulation
bildet \ac{CBS} keine echten Systeme ab, sodass die Übertragbarkeit der
Ergebnisse auf reale Netze offenbleibt. Die höhere Wirklichkeitstreue der
emulationsbasierten Ansätze wäre hier gegen die Zahl der durchführbaren Läufe
einzutauschen gewesen. Da die Fragestellung auf den \emph{Vergleich} von
Topologien unter sonst gleichen Bedingungen zielt und nicht auf absolute
Aussagen über konkrete Netze, wiegt die Zahl der Läufe schwerer.
\Cref{sec:limitationen} greift diesen Punkt erneut auf.

\subsection{Das Lernverfahren}

MARLon stellt mit \ac{A2C} und \ac{PPO} zwei Verfahren aus \ac{SB3} bereit
(\Cref{sec:cbs}). Für beide existieren vollständige, parallel gepflegte
Implementierungen für Angreifer wie Verteidiger, sodass die Wahl offensteht.
Sie fiel auf \ac{PPO}.

Die in \Cref{sec:ppo} beschriebene Deckelung erzeugt gleichmäßigere
Lernkurven, da einzelne Ausreißer den bisherigen Lernfortschritt nicht in einem
Schritt verwerfen können. Betroffen ist davon nicht allein die Darstellung,
sondern die Messbarkeit: Zu den erhobenen Größen gehört die Lerneffizienz und
damit der Zeitpunkt, zu dem ein Lauf als konvergiert gilt, und
das Konvergenzkriterium in \Cref{sec:konvergenz} entscheidet anhand der
Streuung und des Trends der Episoden-Rewards innerhalb eines Fensters. Ein
Verfahren, dessen Kurve wiederholt springt, würde dieses Kriterium unabhängig
von der Topologie später auslösen und die Vergleichbarkeit zwischen den
Topologien beeinträchtigen. Die stabilere Kurve ist somit Voraussetzung dafür,
dass die gemessenen Unterschiede der Topologie zugerechnet werden können.

Für den Diskontfaktor aus \Cref{sec:mdp} gilt in dieser Arbeit
$\gamma = 0{,}99$; ein Reward verliert pro Zeitschritt somit ein Prozent an
Gewicht. Der Wert wurde nicht eigens gewählt, sondern entspricht dem
Standardwert von \ac{SB3}~\cite{stablebaselines3} und dem in der
\ac{PPO}-Literatur gebräuchlichen Ansatz~\cite{schulman2017ppo}. Die übrigen
Hyperparameter sind in \Cref{sec:hyperparameter} aufgeführt.

\section{Umgebung und Kontrolle der Variablen}
\label{sec:umgebung}

Jede Topologie besteht aus denselben zehn Knoten mit denselben Werten. Der
Einstiegsknoten \texttt{WebServer} trägt den Wert~0, die Arbeitsplatzrechner
je~10, der Mail- und der Applikationsserver je~15, der Dateiserver~20, der
Backup-Server~25, der Domain Controller~30 und der Datenbankserver als Crown
Jewel~100. Die Summe dieser Knotenwerte beträgt in jeder Topologie exakt
235~Punkte.

Das Crown Jewel im Sinne von \Cref{sec:angreifer-verteidiger} ist hier der
Datenbankserver. Bei ihm laufen die Daten zusammen, auf denen der
Geschäftsbetrieb beruht, und er vereint sie zugleich an einer einzigen Stelle;
gerade diese Bündelung macht ihn zum lohnenden Ziel, denn ein einzelner
erfolgreicher Zugriff genügt, um den gesamten Datenbestand zu entwenden. Alle
übrigen Systeme sind für einen Angreifer demgegenüber vor allem als
Zwischenstationen von Interesse. Angenommen wird damit ein Angriff, der auf
das Entwenden von Daten zielt, die \emph{Exfiltration}; bei einem anderen
Angreiferziel fiele die Wahl möglicherweise anders aus.

Dass jedem System überhaupt ein einzelner Zahlenwert zugeordnet wird, ist keine
Eigenheit dieser Simulation, sondern folgt dem Vorgehen der in
\Cref{sec:angreifer-verteidiger} beschriebenen Crown Jewels Analysis: Dort ist
die Zuweisung eines Relevanzwerts je IT-System der Kern des Verfahrens, und aus
diesen Werten ergibt sich anschließend, welche Systeme als besonders
schützenswert gelten~\cite[Abschn.~II.A]{gangupantulu2021cja}. Die hier
vergebenen Werte bilden diese Rangfolge nach: Die Produktionsdatenbank trägt
mit 100~Punkten ein Vielfaches jedes anderen Systems, der Verzeichnisdienst
folgt mit 30~Punkten als wichtigster Wegbereiter, und die Arbeitsplatzrechner
bilden mit je 10~Punkten das untere Ende. Der Einstiegsknoten trägt den
Wert~0, weil sein Besitz dem Angreifer nach dem Assumed-Breach-Prinzip von
Beginn an zusteht und deshalb keinen Fortschritt darstellt.

Diese Gleichheit ist beabsichtigt: Da Knotenzahl und Knotenwerte identisch
bleiben, unterscheiden sich die Topologien ausschließlich darin, welche Knoten
miteinander kommunizieren dürfen und wo die Zugangsdaten liegen. Ein
Unterschied im Lernverhalten lässt sich somit auf die Struktur zurückführen und
nicht auf ein lukrativeres Ziel.

Die Knotenwerte sind allerdings nicht der einzige Beitrag zum Reward.
\Cref{sec:rewards} führt die vollständige Bewertung beider Agenten auf und
zeigt, dass die erreichbaren Reward-Niveaus zwischen den Topologien erheblich
auseinanderliegen.

Als Verteidigungsumgebung dienen die vier in \Cref{sec:segmentierungsmuster}
qualitativ beschriebenen Muster, umgesetzt als \emph{flat},
\emph{hub\_and\_spoke}, \emph{dmz} und \emph{micro\_segmented}. Hinzu kommt
\emph{chain}: Diese Topologie dient ausschließlich dem Angreifertraining und
wird nicht verteidigt.

\subsection{Die Topologien als Zehn-Knoten-Modell}
\label{sec:topomodell}

Angestrebt ist ein Netz, das ein Leser als Unternehmensnetz wiedererkennt: ein
von außen erreichbarer Webserver, dahinter Anwendungs-, Mail- und
Dateidienste, ein Verzeichnisdienst und eine Produktionsdatenbank als
wertvollstes System. Die Rollen und ihre Beziehungen sind so gewählt, dass die
Wege zwischen ihnen fachlich begründbar bleiben und nicht bloß eine abstrakte
Graphenstruktur abbilden.

Diesem Anspruch steht die Vergleichbarkeit gegenüber, und wo beide in Konflikt
geraten, hat die Vergleichbarkeit Vorrang. Weil dasselbe Inventar aus zehn
Knoten in allen Topologien gelten muss, lassen sich keine Komponenten
ergänzen, die nur in einer davon vorkämen. Am deutlichsten zeigt sich das beim
Hub-and-Spoke-Modell: Wie in \Cref{sec:segmentierungsmuster} beschrieben, wäre
der Hub in einem realen Netz ein eigens für diese Aufgabe vorgesehenes und
entsprechend gehärtetes Gerät, das ein Angreifer nicht als Sprungbrett nutzt.
Ein solcher Knotentyp existiert im Modell nicht, weshalb die Rolle einem der
vorhandenen Systeme zufällt, nämlich dem Applikationsserver; die Wahl begründet
die Beschreibung der Topologie weiter unten. Diese Vereinfachung ist bewusst in
Kauf genommen; ihre Folgen für die Aussagekraft greift \Cref{sec:limitationen}
auf.

Zum Verständnis der folgenden Abbildungen sind zwei Kantenarten zu
unterscheiden, weil das Modell beides bewusst trennt. Eine durchgezogene Kante
bedeutet \emph{Erreichbarkeit}: Die Firewall lässt Verkehr zum Zielknoten zu.
Eine gestrichelte Kante bedeutet \emph{Besitz von Zugangsdaten}: Auf dem
Ausgangsknoten sind die Anmeldedaten des Zielknotens hinterlegt.

Beide Kantenarten sind dabei gerichtet und beschreiben ausschließlich, wohin
sich ein Angreifer bewegen kann. Sie bilden nicht den gesamten Verkehr eines
Netzes ab. In einem realen Unternehmensnetz greift ein Arbeitsplatzrechner
selbstverständlich auf Server zu; solche Verbindungen sind hier nur dort
verzeichnet, wo sie einem Angreifer einen Weg eröffnen. Ein Knoten ohne
eingehende Kante ist deshalb nicht vom Netz abgeschnitten, sondern für den
Angreifer schlicht kein mögliches nächstes Ziel.

Beide Bedingungen müssen erfüllt sein, damit ein Knoten übernommen werden kann,
sie müssen aber nicht von derselben Stelle herrühren. Erbeutete Zugangsdaten
bleiben dem Angreifer erhalten und lassen sich auch von einem anderen Knoten aus
einsetzen, sofern von dort ein Pfad zum Ziel führt. Erreichbarkeit und
Berechtigung können also aus zwei verschiedenen Quellen stammen. Welche
Verbindungsversuche der Angreifer darüber hinaus überhaupt unternehmen kann,
schränkt \Cref{sec:mask-reichweite} weiter ein.

Eine Besonderheit gilt für den Domain-Controller: Er hält in jeder Topologie
die Zugangsdaten sämtlicher Knoten und bildet damit die in
\Cref{sec:angreifer-verteidiger} beschriebene Rolle eines
Active-Directory-Servers nach. Wer ihn übernimmt, verfügt über die Zugangsdaten
zum gesamten Netz; ob er die betreffenden Systeme auch erreicht, entscheidet
weiterhin die Topologie. Die acht zugehörigen Kanten sind in den Abbildungen
aus Gründen der Lesbarkeit durch eine Beschriftung ersetzt.

\medskip\noindent\topolegende

\paragraph{flat.}
Ein flaches Netz kennt keine interne Segmentierung; jeder Knoten erreicht jeden
anderen (\Cref{fig:topo-flat}). Die fehlende Segmentierung betrifft dabei
ausschließlich die Erreichbarkeit. Die Zugangsdaten sind auch hier nicht
beliebig verteilt, sondern folgen der tatsächlichen Nutzung der Systeme. Alle
drei Arbeitsplätze binden dieselbe Dateiablage ein, denn genau dafür existiert
ein Dateiserver. Der Webserver spricht als Darstellungsschicht nur mit seinem
Anwendungs-Backend, und dieses hält seinerseits einen fest konfigurierten
Datenbankzugang sowie die Anmeldedaten für den Verzeichnisdienst. Mail- und
Dateiserver melden sich ebenfalls am Verzeichnisdienst an, wie es für
Dienstkonten üblich ist.

Bewusst nicht abgebildet ist der Fall, dass sich ein privilegiertes Konto an
einem Endgerät anmeldet und dort verwertbare Anmeldedaten hinterlässt. Kein
Arbeitsplatz hält Zugangsdaten des Domain-Controllers. Die Modellannahme
lautet, dass Administration von einem gesonderten Gerät außerhalb dieses
Inventars erfolgt. Der Weg zum Verzeichnisdienst führt daher ausschließlich
über die Dienstkonten der Server.

\begin{figure}[H]
  \centering
  % AUTOMATISCH ERZEUGT von gen_topologie_figuren.py -- nicht von Hand aendern.
\begin{adjustbox}{max width=\textwidth}
\begin{tikzpicture}[topofig, scale=1.0]
  \node[entrynode] (WebServer) at (0.0,4.0) {WebServer\\[-1pt]{\tiny 0\,P \textbar\ --}};
  \node[netnode] (Workstation1) at (-3.8,-1.24) {WS1\\[-1pt]{\tiny 10\,P \textbar\ SSH}};
  \node[netnode] (Workstation2) at (-2.35,-3.24) {WS2\\[-1pt]{\tiny 10\,P \textbar\ SSH}};
  \node[netnode] (Workstation3) at (0.0,-4.0) {WS3\\[-1pt]{\tiny 10\,P \textbar\ SSH}};
  \node[netnode] (MailServer) at (2.35,-3.24) {MailServer\\[-1pt]{\tiny 15\,P \textbar\ HTTPS}};
  \node[netnode] (FileServer) at (-3.8,1.24) {FileServer\\[-1pt]{\tiny 20\,P \textbar\ SMB}};
  \node[netnode] (AppServer) at (2.35,3.24) {AppServer\\[-1pt]{\tiny 15\,P \textbar\ HTTP}};
  \node[netnode] (BackupServer) at (-2.35,3.24) {BackupServer\\[-1pt]{\tiny 25\,P \textbar\ SMB}};
  \node[dcnode] (DomainController) at (3.8,-1.24) {DomainCtrl\\[-1pt]{\tiny 30\,P \textbar\ RDP}};
  \node[cjnode] (DatabaseServer) at (3.8,1.24) {DatenbankSrv\\[-1pt]{\tiny 100\,P \textbar\ SQL}};
  % Vollstaendiger Graph: alle Kanten wuerden die Abbildung unlesbar
  % machen; stattdessen Hinweis in der Legende.
  \draw[cred] (WebServer) to[bend right=14] (AppServer);
  \draw[cred] (AppServer) to[bend right=14] (DatabaseServer);
  \draw[cred] (AppServer) to[bend right=14] (DomainController);
  \draw[cred] (MailServer) to[bend right=14] (DomainController);
  \draw[cred] (FileServer) to[bend right=14] (BackupServer);
  \draw[cred] (FileServer) to[bend right=14] (DomainController);
  \draw[cred] (Workstation1) to[bend right=14] (FileServer);
  \draw[cred] (Workstation2) to[bend right=14] (FileServer);
  \draw[cred] (Workstation3) to[bend right=14] (FileServer);
  \node[dcbadge, anchor=north] at (3.8,-1.8599999999999999) {h\"alt alle Zugangsdaten};
\end{tikzpicture}
\end{adjustbox}

  \caption[Topologie \emph{flat}]{Topologie \emph{flat}. Die Erreichbarkeit ist
  vollständig, jeder Knoten erreicht jeden anderen; die 90 gerichteten Kanten
  sind nicht eingezeichnet. Dargestellt sind daher nur die Zugangsdaten.}
  \label{fig:topo-flat}
\end{figure}

\paragraph{dmz.}
Die Zuordnung folgt den drei Zonen (\Cref{fig:topo-dmz}). In der
entmilitarisierten Zone liegen der Webserver sowie die beiden Dienste, die er
benötigt: der Mailserver für ausgehende Benachrichtigungen und der
Applikationsserver als Backend. Beide melden sich am Domain-Controller an und
überschreiten damit als einzige die Zonengrenze ins Intranet. Der
Applikationsserver hält zusätzlich die Zugangsdaten der Datenbank, wie es einer
fest konfigurierten Datenbankverbindung entspricht. Im Intranet selbst greift
keine weitere Trennung mehr: Der Domain-Controller erreicht alle Endsysteme,
und der Dateiserver sichert auf den Backup-Server. Genau dieses Gefälle ist
gemeint, wenn vom befestigten Perimeter bei weichem Inneren die Rede ist.

\begin{figure}[H]
  \centering
  % AUTOMATISCH ERZEUGT von gen_topologie_figuren.py -- nicht von Hand aendern.
\begin{adjustbox}{max width=\textwidth}
\begin{tikzpicture}[topofig, scale=0.95]
  \node[entrynode] (WebServer) at (0.0,3.1) {WebServer\\[-1pt]{\tiny 0\,P \textbar\ --}};
  \node[netnode] (Workstation1) at (10.8,6.0) {WS1\\[-1pt]{\tiny 10\,P \textbar\ SSH}};
  \node[netnode] (Workstation2) at (10.8,4.6) {WS2\\[-1pt]{\tiny 10\,P \textbar\ SSH}};
  \node[netnode] (Workstation3) at (10.8,3.2) {WS3\\[-1pt]{\tiny 10\,P \textbar\ SSH}};
  \node[netnode] (MailServer) at (3.6,4.4) {MailServer\\[-1pt]{\tiny 15\,P \textbar\ HTTPS}};
  \node[netnode] (FileServer) at (10.8,1.8) {FileServer\\[-1pt]{\tiny 20\,P \textbar\ SMB}};
  \node[netnode] (AppServer) at (3.6,1.8) {AppServer\\[-1pt]{\tiny 15\,P \textbar\ HTTP}};
  \node[netnode] (BackupServer) at (10.8,0.4) {BackupServer\\[-1pt]{\tiny 25\,P \textbar\ SMB}};
  \node[dcnode] (DomainController) at (7.2,4.4) {DomainCtrl\\[-1pt]{\tiny 30\,P \textbar\ RDP}};
  \node[cjnode] (DatabaseServer) at (7.2,1.8) {DatenbankSrv\\[-1pt]{\tiny 100\,P \textbar\ SQL}};
  \draw[reach] (WebServer) -- (MailServer);
  \draw[reach] (WebServer) -- (AppServer);
  \draw[reach] (MailServer) -- (DomainController);
  \draw[reach] (AppServer) -- (DatabaseServer);
  \draw[reach] (AppServer) -- (DomainController);
  \draw[reach] (DomainController) -- (Workstation1);
  \draw[reach] (DomainController) -- (Workstation2);
  \draw[reach] (DomainController) -- (Workstation3);
  \draw[reach] (DomainController) -- (FileServer);
  \draw[reach] (DomainController) -- (BackupServer);
  \draw[reach] (DomainController) -- (DatabaseServer);
  \draw[reach] (FileServer) -- (BackupServer);
  \draw[cred] (WebServer) to[bend right=14] (MailServer);
  \draw[cred] (WebServer) to[bend right=14] (AppServer);
  \draw[cred] (MailServer) to[bend right=14] (DomainController);
  \draw[cred] (AppServer) to[bend right=14] (DatabaseServer);
  \draw[cred] (AppServer) to[bend right=14] (DomainController);
  \draw[cred] (FileServer) to[bend right=14] (BackupServer);
  \node[dcbadge, anchor=north] at (7.2,3.7800000000000002) {h\"alt alle Zugangsdaten};
\end{tikzpicture}
\end{adjustbox}

  \caption[Topologie \emph{dmz}]{Topologie \emph{dmz}. Der Einstiegsknoten
  erreicht nur die beiden Dienste der entmilitarisierten Zone; erst über diese
  führt der Weg ins Intranet.}
  \label{fig:topo-dmz}
\end{figure}

\paragraph{hub\_and\_spoke.}
Die Rolle des Hubs übernimmt der Applikationsserver. Nach dem
Domain-Controller unterhält er die meisten fachlich begründeten Verbindungen
über Bereichsgrenzen hinweg: Er liest aus der Datenbank, meldet sich am
Verzeichnisdienst an und greift auf die Dateiablage zu. Zu den übrigen Spokes
besteht keine solche fachliche Beziehung; die Arbeitsplätze etwa nutzen die
Anwendung über den Webserver und nicht unmittelbar. Das Muster verlangt eine
solche Beziehung aber auch nicht: Die Kanten zwischen Hub und Spokes folgen aus
der Transitrolle selbst, die sämtlichen Verkehr über diesen einen Punkt führt,
und fallen bei jeder Wahl des Hubs gleichermaßen an.

Dass nicht der Domain-Controller zum Hub wurde, obwohl er noch mehr Bereiche
verbindet, ist eine bewusste Entscheidung. Dass der Hub überhaupt ein
übernehmbares System sein muss, ist bereits das oben beschriebene
Zugeständnis an das feste Inventar. Der Domain-Controller hält darüber hinaus
die Zugangsdaten sämtlicher Knoten, und beides zusammen machte seine Übernahme
zum Verlust des gesamten Netzes: Der Angreifer erhielte mit einem einzigen
Knoten Reichweite und Zugriff auf alles, und für sechs der neun übrigen Knoten
verkürzte sich der Angriffsweg um einen Schritt, weil ihre Zugangsdaten dann
bereits am Engpass lägen. Die Trennung von Reichweite und Zugriff, die dieses
Muster gerade auszeichnet, wäre aufgehoben.

Allein der Hub besitzt ausgehende Firewall-Regeln, alle übrigen Knoten haben
keine (\Cref{fig:topo-hub}).

Für den Angreifer bedeutet das eine harte Einschränkung. Von einem übernommenen
Endsystem aus kommt er nicht weiter, denn ein Spoke erreicht keinen anderen
Spoke. Jeder Schritt setzt den Hub voraus. Zugleich verschafft ihm der Hub aber
nur Reichweite, nicht Zugriff: Die Zugangsdaten der Arbeitsplätze sowie des
Mail-, Datei- und Backup-Servers liegen ausschließlich beim Domain-Controller.
Der Angreifer muss also erst diesen übernehmen, um die Reichweite des Hubs
überhaupt nutzen zu können.

\begin{figure}[H]
  \centering
  % AUTOMATISCH ERZEUGT von gen_topologie_figuren.py -- nicht von Hand aendern.
\begin{adjustbox}{max width=\textwidth}
\begin{tikzpicture}[topofig, scale=1.15]
  \node[entrynode] (WebServer) at (-6.6,0.0) {WebServer\\[-1pt]{\tiny 0\,P \textbar\ --}};
  \node[netnode] (Workstation1) at (1.22,2.96) {WS1\\[-1pt]{\tiny 10\,P \textbar\ SSH}};
  \node[netnode] (Workstation2) at (-1.22,2.96) {WS2\\[-1pt]{\tiny 10\,P \textbar\ SSH}};
  \node[netnode] (Workstation3) at (-2.96,1.22) {WS3\\[-1pt]{\tiny 10\,P \textbar\ SSH}};
  \node[netnode] (MailServer) at (2.96,1.22) {MailServer\\[-1pt]{\tiny 15\,P \textbar\ HTTPS}};
  \node[netnode] (FileServer) at (-1.22,-2.96) {FileServer\\[-1pt]{\tiny 20\,P \textbar\ SMB}};
  \node[netnode] (AppServer) at (0.0,0.0) {AppServer\\[-1pt]{\tiny 15\,P \textbar\ HTTP}};
  \node[netnode] (BackupServer) at (1.22,-2.96) {BackupServer\\[-1pt]{\tiny 25\,P \textbar\ SMB}};
  \node[dcnode] (DomainController) at (-2.96,-1.22) {DomainCtrl\\[-1pt]{\tiny 30\,P \textbar\ RDP}};
  \node[cjnode] (DatabaseServer) at (2.96,-1.22) {DatenbankSrv\\[-1pt]{\tiny 100\,P \textbar\ SQL}};
  \draw[reach] (WebServer) -- (AppServer);
  \draw[reach] (AppServer) -- (MailServer);
  \draw[reach] (AppServer) -- (DomainController);
  \draw[reach] (AppServer) -- (Workstation1);
  \draw[reach] (AppServer) -- (Workstation2);
  \draw[reach] (AppServer) -- (Workstation3);
  \draw[reach] (AppServer) -- (FileServer);
  \draw[reach] (AppServer) -- (BackupServer);
  \draw[reach] (AppServer) -- (DatabaseServer);
  \draw[cred] (WebServer) to[bend right=14] (AppServer);
  \draw[cred] (AppServer) to[bend right=14] (DatabaseServer);
  \draw[cred] (AppServer) to[bend right=14] (DomainController);
  \node[dcbadge, anchor=north] at (-2.96,-1.8399999999999999) {h\"alt alle Zugangsdaten};
\end{tikzpicture}
\end{adjustbox}

  \caption[Topologie \emph{hub\_and\_spoke}]{Topologie \emph{hub\_and\_spoke}.
  Der Applikationsserver ist als Hub der einzige Übergang zwischen dem
  Einstiegsknoten und allen übrigen Systemen.}
  \label{fig:topo-hub}
\end{figure}

\paragraph{micro\_segmented.}
Für den Einstiegsknoten bleibt genau eine Verbindung übrig: die zum
Applikationsserver, dessen Backend er als Webserver benötigt. Der
Applikationsserver seinerseits darf die Datenbank und den Domain-Controller
erreichen, weil er Anwendungsdaten liest und Anmeldungen prüfen muss. Zu den
Arbeitsplätzen, zum Mail-, Datei- und Backup-Server besteht dagegen keine
fachliche Notwendigkeit, folglich existiert auch keine entsprechende Freigabe.
Diese sechs Knoten sind damit nicht vom Netz getrennt, sie nehmen ihre Arbeit
weiterhin auf und greifen ihrerseits auf Dienste zu; nach dem oben Gesagten ist
solcher Verkehr hier nur nicht verzeichnet, weil er dem Angreifer keinen Weg
eröffnet. Vom Einstiegsknoten aus führt zu ihnen jedoch kein Pfad, sodass für
den Angreifer höchstens 145 der insgesamt 235 Punkte erreichbar bleiben
(\Cref{fig:topo-micro}).

Ein zweiter Unterschied betrifft die Zugangsdaten und wiegt für den Angriffsweg
schwerer. In den übrigen Topologien hinterlegt der Applikationsserver die
Zugangsdaten der Datenbank dauerhaft bei sich, wie es bei einer fest
konfigurierten Datenbankverbindung üblich ist. Hier speichert er sie nicht; sie
liegen ausschließlich beim Domain-Controller. Der Angreifer erreicht die
Datenbank zwar schon vom Applikationsserver aus, kann sie aber erst übernehmen,
nachdem er den Domain-Controller kompromittiert hat. Aus zwei Sprüngen werden
drei.

\begin{figure}[H]
  \centering
  % AUTOMATISCH ERZEUGT von gen_topologie_figuren.py -- nicht von Hand aendern.
\begin{adjustbox}{max width=\textwidth}
\begin{tikzpicture}[topofig, scale=0.88]
  \node[entrynode] (WebServer) at (0.0,3.2) {WebServer\\[-1pt]{\tiny 0\,P \textbar\ --}};
  \node[netnode,isolated] (Workstation1) at (11.1,5.2) {WS1\\[-1pt]{\tiny 10\,P \textbar\ SSH}};
  \node[netnode,isolated] (Workstation2) at (11.1,3.9) {WS2\\[-1pt]{\tiny 10\,P \textbar\ SSH}};
  \node[netnode,isolated] (Workstation3) at (11.1,2.6) {WS3\\[-1pt]{\tiny 10\,P \textbar\ SSH}};
  \node[netnode,isolated] (MailServer) at (11.1,6.5) {MailServer\\[-1pt]{\tiny 15\,P \textbar\ HTTPS}};
  \node[netnode,isolated] (FileServer) at (11.1,1.3) {FileServer\\[-1pt]{\tiny 20\,P \textbar\ SMB}};
  \node[netnode] (AppServer) at (3.6,3.2) {AppServer\\[-1pt]{\tiny 15\,P \textbar\ HTTP}};
  \node[netnode,isolated] (BackupServer) at (11.1,0.0) {BackupServer\\[-1pt]{\tiny 25\,P \textbar\ SMB}};
  \node[dcnode] (DomainController) at (7.2,4.6) {DomainCtrl\\[-1pt]{\tiny 30\,P \textbar\ RDP}};
  \node[cjnode] (DatabaseServer) at (7.2,1.8) {DatenbankSrv\\[-1pt]{\tiny 100\,P \textbar\ SQL}};
  \draw[reach] (WebServer) -- (AppServer);
  \draw[reach] (AppServer) -- (DatabaseServer);
  \draw[reach] (AppServer) -- (DomainController);
  \draw[cred] (WebServer) to[bend right=14] (AppServer);
  \draw[cred] (AppServer) to[bend right=14] (DomainController);
  \node[dcbadge, anchor=north] at (7.2,3.9799999999999995) {h\"alt alle Zugangsdaten};
\end{tikzpicture}
\end{adjustbox}

  \caption[Topologie \emph{micro\_segmented}]{Topologie
  \emph{micro\_segmented}. Blass dargestellte Knoten sind vom Einstiegsknoten
  aus über keinen Pfad erreichbar. Anders als in den übrigen Topologien
  speichert der Applikationsserver die Zugangsdaten der Datenbank
  nicht zwischen; sie liegen ausschließlich beim Domain-Controller.}
  \label{fig:topo-micro}
\end{figure}

\paragraph{chain.}
Die Kette ist als strikte Linie ohne jede Verzweigung ausgelegt
(\Cref{fig:topo-chain}) und damit die einfachste denkbare Anordnung: Wer sie
lösen will, muss nur den Nachbarn entdecken, dessen Zugangsdaten auslesen und
ihn übernehmen, und das neun Mal hintereinander. Ein darauf trainierter
Angreifer beherrscht folglich genau dieses Grundmuster und nichts darüber
hinaus, insbesondere keine Entscheidung zwischen mehreren Wegen.

Damit dient die Kette als Kontrollbedingung für die Frage, ob dieses
Grundmuster allein bereits genügt. Behauptet sich ein solcher Angreifer auch in
den vier verzweigten Topologien, spricht das dafür, dass diese sich für den
Verteidiger weniger unterscheiden, als ihre Struktur nahelegt. Scheitert er
dort, lässt sich umgekehrt annehmen, dass die Verzweigungen eine andere
Fähigkeit verlangen. Ein Beweis ist beides nicht, da nur eine einzelne
Angreiferherkunft betrachtet wird; als Anhaltspunkt taugt der Vergleich
gleichwohl. Voraussetzung dafür ist, dass auch der Kette dasselbe Inventar
zugrunde liegt: dieselben zehn Knoten, dieselben Werte, dieselben
Zugangsports. Verändert ist allein die Anordnung.

\begin{figure}[H]
  \centering
  % AUTOMATISCH ERZEUGT von gen_topologie_figuren.py -- nicht von Hand aendern.
\begin{adjustbox}{max width=\textwidth}
\begin{tikzpicture}[topofig, scale=1.0]
  \node[entrynode] (WebServer) at (0.0,1.6) {WebServer\\[-1pt]{\tiny 0\,P \textbar\ --}};
  \node[netnode] (Workstation1) at (2.4,1.6) {WS1\\[-1pt]{\tiny 10\,P \textbar\ SSH}};
  \node[netnode] (Workstation2) at (4.8,1.6) {WS2\\[-1pt]{\tiny 10\,P \textbar\ SSH}};
  \node[netnode] (Workstation3) at (7.2,1.6) {WS3\\[-1pt]{\tiny 10\,P \textbar\ SSH}};
  \node[netnode] (MailServer) at (9.6,1.6) {MailServer\\[-1pt]{\tiny 15\,P \textbar\ HTTPS}};
  \node[netnode] (FileServer) at (9.6,0.0) {FileServer\\[-1pt]{\tiny 20\,P \textbar\ SMB}};
  \node[netnode] (AppServer) at (7.2,0.0) {AppServer\\[-1pt]{\tiny 15\,P \textbar\ HTTP}};
  \node[netnode] (BackupServer) at (4.8,0.0) {BackupServer\\[-1pt]{\tiny 25\,P \textbar\ SMB}};
  \node[dcnode] (DomainController) at (2.4,0.0) {DomainCtrl\\[-1pt]{\tiny 30\,P \textbar\ RDP}};
  \node[cjnode] (DatabaseServer) at (0.0,0.0) {DatenbankSrv\\[-1pt]{\tiny 100\,P \textbar\ SQL}};
  \draw[reach] (WebServer) -- (Workstation1);
  \draw[reach] (Workstation1) -- (Workstation2);
  \draw[reach] (Workstation2) -- (Workstation3);
  \draw[reach] (Workstation3) -- (MailServer);
  \draw[reach] (MailServer) -- (FileServer);
  \draw[reach] (FileServer) -- (AppServer);
  \draw[reach] (AppServer) -- (BackupServer);
  \draw[reach] (BackupServer) -- (DomainController);
  \draw[reach] (DomainController) -- (DatabaseServer);
  \draw[cred] (WebServer) to[bend right=14] (Workstation1);
  \draw[cred] (Workstation1) to[bend right=14] (Workstation2);
  \draw[cred] (Workstation2) to[bend right=14] (Workstation3);
  \draw[cred] (Workstation3) to[bend right=14] (MailServer);
  \draw[cred] (MailServer) to[bend right=14] (FileServer);
  \draw[cred] (FileServer) to[bend right=14] (AppServer);
  \draw[cred] (AppServer) to[bend right=14] (BackupServer);
  \draw[cred] (BackupServer) to[bend right=14] (DomainController);
  \draw[cred] (DomainController) to[bend right=14] (DatabaseServer);
\end{tikzpicture}
\end{adjustbox}

  \caption[Topologie \emph{chain}]{Topologie \emph{chain}, verwendet für das
  Angreifertraining. Jeder Knoten gibt Erreichbarkeit und Zugangsdaten
  ausschließlich für seinen Nachfolger preis; von dieser Topologie führen neun
  Sprünge bis zum Crown Jewel.}
  \label{fig:topo-chain}
\end{figure}

\section{Trainingsphasen}
\label{sec:phasen}

Der Versuch gliedert sich in zwei Phasen. In der ersten wird ausschließlich der
Angreifer trainiert, in der zweiten ausschließlich der Verteidiger.

\subsection{Phase 1: Angreifer-Vortraining}

Ein Verteidiger lässt sich nur sinnvoll bewerten, wenn ihm ein kompetenter
Gegner gegenübersteht. Deshalb werden zunächst Angreifer ohne jede Verteidigung
trainiert, jeweils \num{500000} Zeitschritte lang.

Naheliegend wäre, den Angreifer schon hier gegen einen Verteidiger antreten zu
lassen, etwa einen zufällig handelnden, um ihn auf spätere Gegenwehr
vorzubereiten. Darauf wird bewusst verzichtet, und zwar aus demselben Grund, aus
dem in Phase~2 der Angreifer eingefroren wird: Was den Angreifer formt, soll
allein die Topologie sein. Ein mittrainierter Verteidiger, ob zufällig oder
lernend, wirkt als zweite Einflussgröße auf die Herkunft jedes Angreifers ein,
und die beobachteten Unterschiede zwischen den Angreifern ließen sich nicht mehr
eindeutig der Struktur zuschreiben. Ein zufälliger Verteidiger wäre zudem selbst
willkürlich gesetzt: Seine Aktionen bevorzugten keine Topologie planvoll,
verzerrten das Angreifertraining aber dennoch in einer nicht kontrollierten
Weise. Der Preis dieser Reinheit ist, dass die Angreifer keine Erfahrung mit
Gegenwehr sammeln; \Cref{sec:limitationen} greift das auf.

Fünf \emph{Spezialisten} entstehen dabei, einer je Topologie. Jeder lernt zwar
nur in einer einzigen Umgebung, tritt in Phase~2 jedoch in allen vier
Verteidigungstopologien als Gegner an. Ein Angreifer trifft dort also in aller
Regel auf eine Struktur, die er im Training nicht gesehen hat. Genau diese
Kreuzung erlaubt es, den Einfluss der verteidigten Topologie von dem der
Angreiferherkunft zu trennen. Ergänzend
wird ein \emph{Generalist} trainiert, im Folgenden Super-Angreifer genannt: Ein
einzelnes Modell durchläuft nacheinander alle fünf Topologien -- in der
Reihenfolge \emph{chain}, \emph{flat}, \emph{hub\_and\_spoke}, \emph{dmz},
\emph{micro\_segmented} --, wobei die Gewichte von Stufe zu Stufe
weitergereicht werden. Damit lässt sich prüfen, ob
ein über mehrere Strukturen trainierter Angreifer den Spezialisten überlegen
ist.

\subsection{Phase 2: Verteidiger-Matrix}

In der zweiten Phase wird für jede Kombination aus Verteidigungstopologie und
Angreifer ein eigener Verteidiger trainiert. Bei vier Topologien und sechs
Angreifern ergeben sich 24~Matchups. Der Angreifer ist dabei \emph{eingefroren}:
Sein Modell wird geladen, handelt weiter, erhält aber keine Gewichtsanpassungen
mehr.

Das Einfrieren hält den Gegner über den gesamten Lauf konstant und macht die
beobachtete Lernkurve dem Verteidiger zurechenbar. Lernten beide Agenten
gleichzeitig, verschöbe sich das Ziel fortlaufend, und ein Unterschied zwischen
zwei Topologien ließe sich nicht mehr von einem Unterschied im Gegnerverhalten
trennen. Diese Entscheidung lässt sich auch am MARLon-Aufbau selbst
festmachen: Dort führt gleichzeitiges Training beider Agenten dazu, dass der
Angreifer keinen brauchbaren Angriffsweg mehr findet, weil wiederholte
Rücksetzungen durch einen anfangs instabilen Verteidiger ihn immer wieder auf
den Ausgangszustand zurückwerfen, bevor er über den Anfang der Episode
hinauskommt~\autocite{marlon}. Erkauft wird die Zurechenbarkeit damit, dass der
Verteidiger sich auf genau ein unveränderliches Gegnerverhalten einspielen kann.
Eben deshalb wird nicht ein einzelner Angreifer eingesetzt, sondern gegen jede
Topologie alle sechs: Jeder Verteidiger sieht zwar nur einen Gegner, über die
Matrix hinweg werden aber sechs verschiedene Angriffsmuster geprüft, davon fünf,
die nicht auf der verteidigten Topologie entstanden sind. Vollständig aufheben
lässt sich die Einschränkung dadurch nicht; \Cref{sec:limitationen} greift sie auf.

Jedes Matchup wird fünfmal wiederholt, und verglichen wird nicht der einzelne
Lauf, sondern der über die fünf Wiederholungen gemittelte Verlauf. Insgesamt
ergeben sich damit \num{120}~Trainingsläufe zu je \num{500000} Zeitschritten.

\section{Aktionsraum und Invalid-Action-Masking}
\label{sec:masking}

Was ein Agent lernen kann, ist durch die Menge der ihm zur Verfügung
stehenden Aktionen begrenzt. Dieser Abschnitt beschreibt diese Menge für
beide Seiten samt der Maskierung, die im jeweiligen Zustand ungültige
Aktionen ausschließt, und prüft abschließend, welche mehrstufigen
Strategien darin überhaupt angelegt sind.

\subsection{Der Aktionsraum des Angreifers}
\label{sec:aktionsraum-atk}

In jedem Zeitschritt wählt der Angreifer genau eine Aktion aus. Eine Aktion ist
dabei keine grobe Absicht wie \glqq greife den Dateiserver an\grqq{}, sondern eine
vollständig ausformulierte Anweisung: Sie legt fest, was getan wird, von
welchem Knoten aus und gegen welches Ziel. \ac{CBS} kennt dafür drei Arten.

Eine \emph{lokale Aktion} durchsucht einen bereits übernommenen Knoten von
innen, etwa nach dort gespeicherten Zugangsdaten. Eine \emph{entfernte Aktion}
richtet sich von einem übernommenen Knoten aus gegen ein anderes System und
späht dieses aus. Eine \emph{Verbindung} meldet sich mit erbeuteten
Zugangsdaten an einem Zielknoten an und übernimmt ihn. Nur die dritte Art
erweitert den Besitz; die ersten beiden beschaffen die Voraussetzungen dafür.

Jede Kombination aus Aktionsart und Parametern bildet einen eigenen Eintrag.
Die Größe des Aktionsraums ergibt sich damit als Produkt der jeweiligen
Auswahlmöglichkeiten (\Cref{tab:aktionsraum}).

\begin{table}[H]
  \centering
  \caption[Aufbau des Aktionsraums]{Aufbau des Aktionsraums des Angreifers bei
  zehn Knoten.}
  \label{tab:aktionsraum}
  \begin{tabular}{llr}
    \toprule
    Art & Parameter & Anzahl \\
    \midrule
    Lokale Aktion   & 10 Knoten $\times$ 2 Schwachstellen & 20 \\
    Entfernte Aktion & 10 Quellen $\times$ 10 Ziele $\times$ 1 Schwachstelle & 100 \\
    Verbindung      & 10 Quellen $\times$ 10 Ziele $\times$ 6 Ports $\times$ 9 Zugangsdaten & 5400 \\
    \midrule
    Gesamt          &  & 5520 \\
    \bottomrule
  \end{tabular}
\end{table}

Die 5400 Verbindungen wirken auf den ersten Blick redundant, denn ein Knoten
besitzt in der Modellierung genau einen Zugangsport und genau einen Satz
Zugangsdaten; mit der Wahl des Ziels stünden beide inhaltlich also bereits
fest. \ac{CBS} nimmt diese Einschränkung jedoch nicht vor:
Angeboten werden für jede Quelle-Ziel-Kombination sämtliche 54 Kombinationen
aus Port und zwischengespeicherter Zugangsdatenkarte, mit dem ausdrücklichen
Hinweis im Quellcode der Umgebung, dass Erfolg damit nicht garantiert sei. Der
Angreifer muss also selbst lernen, welche einzelne davon zum jeweiligen Knoten
passt; sie sind dem Aktionsraum nicht vorweg entzogen.

Technisch sind diese 5520 Einträge als flaches \texttt{Discrete} umgesetzt,
also als eine einzige durchnummerierte Liste, in der jede vollständige Aktion
einen eigenen Index besitzt. MARLon verwendet an dieser Stelle ein
\texttt{MultiDiscrete}, bei dem der Agent je Zug eine Reihe von Zahlen ausgibt:
eine für die Aktionsart, weitere für Quellknoten, Zielknoten, Port und
Zugangsdaten.

Nötig wurde die Umstellung durch das im Folgenden beschriebene Masking. Dieses
kennzeichnet jeden Eintrag der Liste einzeln als gültig oder ungültig. Bei einem
\texttt{MultiDiscrete} müsste jede Stelle für sich gekennzeichnet werden, und
damit lassen sich verknüpfte Bedingungen nicht ausdrücken: Ob ein Zielknoten in
Frage kommt, hängt davon ab, welcher Quellknoten gewählt wurde, denn nur von
einem übernommenen Knoten aus lässt sich überhaupt handeln. Eine Kennzeichnung
je Stelle kennt diese Abhängigkeit nicht und müsste einen Zielknoten entweder
immer oder nie zulassen. Erst wenn sämtliche Kombinationen als einzelne
Einträge nebeneinanderstehen, lässt sich die Bedingung abbilden, weil dann genau
die gültigen \emph{Gesamt}aktionen bezeichnet werden können. Der Preis ist eine
sehr lange Liste, doch da davon je Schritt nur wenige Einträge gültig sind,
wählt der Angreifer tatsächlich aus einer kleinen Menge.

Das Verfahren selbst ist das \emph{Invalid-Action-Masking} nach Huang und
Ontañón~\autocite{huang2022masking}. Vor jeder Entscheidung wird
ein boolescher Vektor übergeben, der die im aktuellen Zustand gültigen Aktionen
markiert; die Policy wählt ausschließlich aus diesen. Ohne Maskierung entfällt
ein Großteil der Schritte auf wirkungslose Aktionen, etwa den Angriff von einem
nicht kontrollierten Knoten aus.

Technisch gibt der Actor je Aktion zunächst einen unnormierten Zahlenwert aus,
ein sogenanntes \emph{Logit}; erst ein nachgelagerter Rechenschritt überführt
diese Werte in die Wahrscheinlichkeiten der
Policy~\cite[Abschn.~\glqq Background\grqq{}]{huang2022masking}. Die Maskierung
setzt an dieser Stelle an: Die Logits ungültiger Aktionen werden vor der
Normierung durch eine sehr große negative Zahl ersetzt, in der hier
eingesetzten Implementierung $-10^{8}$, sodass ihre Wahrscheinlichkeit
praktisch null wird. Die Wahrscheinlichkeitsmasse verteilt sich dadurch
ausschließlich auf die gültigen Aktionen, statt über alle 5520 zu streuen. Die
Markov-Eigenschaft bleibt davon unberührt, denn maskiert wird nicht der Zustand,
sondern die in ihm zulässige Aktionsmenge. Zustandsabhängige Aktionsmengen sind
Bestandteil der üblichen \ac{MDP}-Definition. Huang und Ontañón zeigen zudem,
dass der resultierende Gradient weiterhin ein gültiger Policy-Gradient
ist~\autocite{huang2022masking}.

Entscheidend ist der Zeitpunkt der Maskenberechnung. Sie erfolgt unmittelbar
vor der Aktionswahl aus dem aktuellen Zustand. Eine zwischengespeicherte Maske
wäre um den vorangegangenen Verteidigerzug veraltet, da dieser Knoten neu
aufsetzen oder Verbindungen sperren kann. Huang und Ontañón behandeln den Fall
eines einzelnen Agenten, bei dem diese Frage nicht auftritt; sobald ein zweiter
Agent zwischen zwei Zügen handelt, genügt es nicht mehr, die Maske einmal je
eigenem Schritt zu bestimmen.

\subsubsection{Reichweite der Maskierung: Zugangsdaten allein genügen nicht}
\label{sec:mask-reichweite}

Die von \ac{CBS} bereitgestellte Aktionsmaske prüft für eine Verbindung
ausschließlich, ob der Angreifer die betreffende Zugangsdatenkarte besitzt,
nicht aber, ob zwischen Quelle und Ziel überhaupt ein Netzpfad besteht.
Weil die Modellierung Erreichbarkeit und Zugangsdatenbesitz trennt
(\Cref{sec:topomodell}), würde sie damit Verbindungen entlang gar nicht
vorhandener Kanten erlauben, sobald der Angreifer die passenden
Zugangsdaten irgendwo erbeutet hat, und die Topologie als Einflussgröße
unterlaufen.

Die native Maske wurde deshalb um eine eigene Einschränkung ergänzt. Der Angreifer führt seither mit, welche Kante er von welchem Knoten aus
tatsächlich aufgedeckt hat; aufgedeckt wird eine Kante ausschließlich durch
die auf den Quellknoten gerichtete lokale Aktion zur
Nachbarschaftserkundung, nicht durch das Erbeuten von Zugangsdaten. Die von
\ac{CBS} gelieferte Maske wird für Verbindungen und entfernte Aktionen mit
dieser Liste geschnitten und bildet damit auch die Netzstruktur ab.

Anders als das Masking selbst, das bereits im Referenzlauf eingesetzt wurde,
stammt diese Erweiterung aus der späteren Überarbeitung. Sie greift unmittelbar
in das Verhalten des Angreifers ein, weshalb sämtliche Angreifer für die
berichteten Läufe neu vortrainiert wurden.

\subsection{Der Aktionsraum des Verteidigers}
\label{sec:aktionsraum-def}

Der Verteidiger verfügt über drei Aktionsarten: Er kann einen Knoten \emph{neu
aufsetzen} und damit den Angreifer von diesem System vertreiben, er kann
dessen eingehenden \emph{Verkehr sperren}, und er kann eine bestehende Sperre
wieder \emph{aufheben}. Jede der drei Arten benennt ausschließlich den
betroffenen Knoten; Sperren und Aufheben wirken dabei auf sämtliche
eingehenden Dienste dieses Knotens zugleich, ohne dass Port oder Richtung
gesondert gewählt würden.

Technisch ist das als \texttt{MultiDiscrete} über vier Dimensionen umgesetzt,
$[3, 10, 10, 10]$. Der erste Wert wählt die Aktionsart, die übrigen drei
benennen den Knoten für Neuaufsetzen, Sperren beziehungsweise Aufheben; nur
die zur gewählten Art passende Dimension wird ausgewertet
(\Cref{tab:aktionsraum-def}).

\begin{table}[H]
  \centering
  \caption[Aktionsraum des Verteidigers]{Aktionsraum des Verteidigers bei zehn
  Knoten.}
  \label{tab:aktionsraum-def}
  \begin{tabular}{llrr}
    \toprule
    Aktionsart & Parameter & Kombinationen & davon wirksam \\
    \midrule
    Knoten neu aufsetzen & 10 Knoten & 10 & 9 \\
    Verkehr sperren      & 10 Knoten & 10 & 9 \\
    Verkehr aufheben     & 10 Knoten & 10 & 9 \\
    \midrule
    Gesamt                &  & 30 & 27 \\
    \bottomrule
  \end{tabular}
\end{table}

Der Einstiegsknoten schränkt beide Spalten der Tabelle zusätzlich ein. Er ist
nicht neu aufsetzbar, wodurch für das Neuaufsetzen nur neun der zehn Knoten
tatsächlich gültige Ziele sind. Eine Sperre auf ihm ist zwar technisch gültig,
bringt aber nie einen Nutzen: Er trägt den Wert null und bleibt über die
gesamte Episode kompromittiert, weil er nicht zurückgewonnen werden kann.
Wirksam wählbar bleiben damit 27 der 30 strukturell möglichen Aktionen.

Gegenüber dem Original ist dieser Aktionsraum damit in zweifacher Hinsicht
enger. MARLon sieht fünf Aktionsarten vor, neben dem Neuaufsetzen, dem Sperren
und dem Aufheben auch das Stoppen und das Starten eines Dienstes, und verlangt
zu jeder von ihnen zusätzlich Port und Richtung als selbst zu wählende
Parameter (\Cref{sec:marlon-original}). Hier sind es drei Aktionsarten mit dem
Knoten als einzigem Parameter. Beide Änderungen gehen auf Beobachtungen am
Referenzlauf zurück.

Als erste Änderung sind Dienst stoppen und Dienst starten entfallen. Für den Angreifer ist
ein gestoppter Dienst von einer gesperrten Verbindung ohnehin nicht zu
unterscheiden, beide lassen den Verbindungsversuch scheitern. Da eine Sperre
nach dem Umbau denselben Knoten in seiner Gesamtheit betrifft und seit der in
\Cref{sec:sla} beschriebenen Korrektur ebenfalls Verfügbarkeit kostet, blieb
kein Zweck, den nicht auch sie erfüllt.

Ausschlaggebend war damit nicht, dass eine der beiden Aktionen die schlechtere
gewesen wäre, sondern dass zwei Aktionen für dieselbe Wirkung den Aktionsraum
ohne Gewinn an Ausdruckskraft vergrößert hätten. Dass die Wahl auf die Sperre
fiel, hat einen inhaltlichen Grund: Sie setzt am Zugang zu einem System an und
lässt dessen Betrieb unberührt, während ein Dienststopp in das System selbst
eingreift. Von zwei gleich wirkenden Mitteln ist damit das mildere gewählt.

Offenzulegen ist dabei, dass diese Gleichwertigkeit nur für die Sperre in ihrer
heutigen Form gilt und nie gemessen wurde. Im Referenzlauf standen dem
Verteidiger zwar beide Aktionen gleichzeitig zur Verfügung, aber nicht als
gleichwertige Alternativen: Die Sperre wirkte dort je Port statt je Knoten,
erreichte wegen der in \Cref{sec:marlon-fehler} beschriebenen Portliste die
wertvollsten Systeme überhaupt nicht und brachte, da sie weder Verfügbarkeit
kostete noch belohnt wurde, keinerlei Signal hervor. Der Dienststopp dagegen
war auf sauberen Knoten mit zehn Punkten belastet. Dass der Verteidiger ihn
dort ablegte, ist deshalb kein Vergleichsergebnis. Ein Lauf, in dem beide
Aktionen unter den heutigen Bedingungen nebeneinander zur Wahl stünden, wurde
nicht durchgeführt; die Gleichwertigkeit folgt aus der Betrachtung der
Wirkungsketten, nicht aus einer Messung.

Als zweite Änderung wählt der Verteidiger für Sperren und Freigaben nicht mehr selbst,
welcher Port betroffen ist. Er benennt nur den Knoten; welche Dienste dort
laufen, ist als Grundwissen des Betreibers vorausgesetzt und kein
Erkundungsergebnis, wie es für den Angreifer eines wäre. Dieselbe Asymmetrie
gilt bereits an anderer Stelle des Modells: Der Verteidiger spricht jeden
Knoten direkt über seinen Index an und darf ihn neu aufsetzen, sieht sein
eigenes Netz für die Zwecke dieser Aktionen also vollständig. Die Wahl des
Ports blieb die einzige Stelle, an der er über sein eigenes Netz raten musste;
diese Uneinheitlichkeit ist mit der Änderung beseitigt.

\subsection{Mehrstufige Strategien im Aktionsraum}
\label{sec:mehrstufig}

Die Wirkung einer Verteidigeraktion hängt davon ab, wo sie ansetzt. Eine
Sperre wirkt nur auf den \emph{Ziel}knoten einer Verbindung. Auf einem Knoten,
den der Angreifer bereits besitzt, ändert eine Sperre deshalb nichts, denn als
Ausgangspunkt einer weiteren Verbindung braucht er den eigenen eingehenden
Zugang nicht. Auf einem noch nicht übernommenen Knoten verhindert sie die
Übernahme über diesen Weg sehr wohl.

Daraus folgt eine Strategie, die den Angreifer vollständig einschließen würde:
Man drängt ihn durch wiederholtes Neuaufsetzen auf den Einstiegsknoten zurück
und sperrt anschließend alle von dort aus erreichbaren Knoten. Der Angreifer
besäße dann zwar weiterhin ein System, käme aber nirgendwo mehr hin. Der
Aktionsraum sieht diese Möglichkeit demnach vor.

Ob ein lernendes Verfahren sie auch findet, ist eine andere Frage. Die
Einschließung besteht aus vielen einzelnen Sperren, und ihr Nutzen tritt erst
ein, wenn die letzte von ihnen gesetzt ist: Solange auch nur ein Weg offen
bleibt, nimmt der Angreifer ihn. Bis dahin bringt jede einzelne Sperre für sich
genommen nichts. Der Verteidiger müsste eine lange Folge bestimmter Aktionen
zufällig treffen, bevor sich überhaupt ein Reward zeigt, an dem \ac{PPO} sie
erkennen könnte. Das ist der in \Cref{sec:credit} beschriebene Fall der
verzögerten Belohnung in ausgeprägter Form.

Erreichbar wird eine solche Strategie erst, wenn bereits die einzelne Sperre
einen eigenen Reward trägt, wenn die Zwischenschritte also für sich verstärkt
werden und der Agent sie nicht erst am Ende der Kette bewertet bekommt. Genau
darauf zielt der in \Cref{sec:abwehrbonus} beschriebene Abwehrbonus. Ob er
dafür ausreicht, wird in \Cref{ch:ergebnisse} untersucht.

\section{Reward-Struktur und Randbedingungen}
\label{sec:rewards}

Der Reward entsteht auf zwei Ebenen: \ac{CBS} bewertet jede Angreiferaktion,
und der Wrapper setzt darauf die Kopplung beider Agenten sowie die in diesem
Abschnitt beschriebenen Ergänzungen.

\begin{longtable}{EWK}
    \caption[Reward des Angreifers]{Reward des Angreifers je Aktion. Aufgeführt
    sind nur die Bestandteile, die in den hier verwendeten Topologien überhaupt
    auftreten können. Der untere Block stammt nicht aus \ac{CBS}, sondern wird
    vom Wrapper ergänzt.}
    \label{tab:reward-atk} \\
    \toprule
    Ereignis & Reward & Anmerkung \\
    \midrule
    \endfirsthead
    \toprule
    Ereignis & Reward & Anmerkung \\
    \midrule
    \endhead
    \midrule
    \multicolumn{3}{@{}r@{}}{\footnotesize Fortsetzung auf der nächsten Seite} \\
    \endfoot
    \bottomrule
    \endlastfoot
    Knoten übernommen             & $+$Knotenwert & 10 bis 100, auch nach Rückeroberung \\
    Knoten neu entdeckt           & $+5$  & je Knoten \\
    Zugangsdaten neu entdeckt     & $+3$  & je Credential \\
    Eigenschaft neu entdeckt      & $+2$  & je Eigenschaft eines Knotens, hier Betriebssystem und Rolle, die ein Scan aus der Ferne sichtbar macht \\
    Aktion erstmals erfolgreich   & $+7$  & je Knoten und Schwachstelle \\
    Fehlversuch, Wiederholung, Aktionskosten & $0$ & in \ac{CBS} zwischen $-1$ und $-50$, durch die unten beschriebene Kappung ohne Wirkung \\
    Alle Knoten übernommen        & $+5000$ & ersetzt den Schrittwert \\
    \midrule
    Knoten gehalten               & $+0{,}005\cdot$Knotenwert & je Knoten und Schritt \\
    Knoten durch Reimage verloren & $-$Knotenwert & \\
\end{longtable}

Damit ist der Reward des Angreifers vollständig beschrieben. Der des
Verteidigers folgt in \Cref{tab:reward-def}. Er entsteht überwiegend durch
Spiegelung der eben aufgeführten Werte und trägt darüber hinaus einige eigene
Posten.

\begin{longtable}{EWK}
    \caption[Reward des Verteidigers]{Reward des Verteidigers je Schritt.}
    \label{tab:reward-def} \\
    \toprule
    Ereignis & Reward & Anmerkung \\
    \midrule
    \endfirsthead
    \toprule
    Ereignis & Reward & Anmerkung \\
    \midrule
    \endhead
    \midrule
    \multicolumn{3}{@{}r@{}}{\footnotesize Fortsetzung auf der nächsten Seite} \\
    \endfoot
    \bottomrule
    \endlastfoot
    Grundkopplung             & $-$Angreifer-Reward & jeder Schritt \\
    Abwehrbonus & $\pm 0{,}00025\cdot$Knotenwert & je Knoten, der gesperrt und zugleich nicht kompromittiert ist; Vorzeichen je nach Bedrohungslage, \Cref{sec:abwehrbonus} \\
    Ungültige Aktion          & $0$ & Zug wird übersprungen \\
    \midrule
    \ac{SLA} unterschritten   & $-50$ & je Schritt, ergänzt den Schrittwert \\
    Angreifer vollständig entfernt & $+5000$ & strukturell unerreichbar \\
\end{longtable}

Mehrere Eigenschaften dieser Struktur sind für die Auswertung wesentlich und
nicht offensichtlich.

Erstens erklärt sich die Nullzeile in \Cref{tab:reward-atk} aus einer Kappung:
\ac{CBS} kennt für Fehlversuche, Wiederholungen und Aktionskosten durchaus
Abzüge, begrenzt den Schrittwert des Angreifers aber nach unten auf null, bevor
er zurückgegeben wird. Ein Fehlversuch kostet ihn also nichts, er bringt
lediglich nichts ein. Diese Kappung ist für sich genommen
sinnvoll: Sie verhindert, dass der Verteidiger über die Grundkopplung Reward
für Vorgänge gutgeschrieben bekäme, zu denen er nichts beigetragen hat, was
das Lernsignal verfälschen würde. Sie schließt aber zugleich, in Verbindung mit
der reinen Grundkopplung, den Kanal, über den Prävention im Reward sichtbar
würde, wie \Cref{sec:marlon-original} gezeigt hat. Die Begrenzung greift
zudem nur auf dem Wert, den \ac{CBS} selbst berechnet; der untere
Tabellenblock wird erst danach aufgeschlagen. Negative Angreifer-Rewards
entstehen deshalb ausschließlich durch verlorene Knoten: Setzt der
Verteidiger einen übernommenen Knoten neu auf, verliert der Angreifer dessen
Wert.

Zweitens folgt aus dem Abzug für verlorene Knoten eine Bedingung, die \ac{CBS}
von sich aus nicht erfüllt. Der Abzug ist eine eigene Ergänzung
(\Cref{tab:reward-atk}, unterer Block): Im Auslieferungszustand kostet ein
Reimage den Angreifer nichts, er behält den einmal gutgeschriebenen Knotenwert
einfach. Damit ist dort auch nichts unausgeglichen. Sobald ein Verlust aber
etwas kostet, muss die Rückeroberung desselben Knotens ihn wieder einbringen,
sonst ist der Kreislauf aus Verlust und Rückeroberung kein Nullsummengeschäft
mehr.

Genau hier lag die Ursache des in \Cref{sec:referenzlauf} geschilderten
Verhaltens, das dem Referenzlauf vorausging. \ac{CBS} schreibt einen Knotenwert
nur beim \emph{ersten} Besitz gut und vermerkt die Übernahme dauerhaft; ein
Reimage löscht diesen Vermerk nicht. Zusammen mit dem neu eingeführten Abzug
erhielt der Angreifer für eine Rückeroberung also nichts, während ihn jeder
Verlust den vollen Knotenwert kostete. Über die Grundkopplung wurde daraus eine
Einnahmequelle für den Verteidiger: Er konnte denselben Knoten wieder und
wieder neu aufsetzen und dabei jedes Mal dessen Wert vereinnahmen, ohne dass
sich an der Lage im Netz etwas änderte. In Vorversuchen trat dieses Verhalten
deutlich hervor. Bei zehn vorhandenen Knoten wurden bis zu 268 Rückgewinnungen
je Episode gezählt, und der Verteidiger-Reward korrelierte mit der Zahl der
Reimages bis zu $r = +0{,}98$. Der Vermerk wird beim Verlust eines Knotens
deshalb zurückgesetzt. Jede Rückeroberung zählt damit erneut als erste und
gleicht den vorangegangenen Abzug exakt aus. Der Verteidiger behält seinen
Gewinn nur dann, wenn der Angreifer den Knoten nicht zurückholt.

Drittens maß der Reward vor der hier vorgenommenen Ergänzung nur, \emph{ob}
ein Knoten kompromittiert wurde, nicht \emph{wie lange}. Ein Angreifer, der ein
System zwei Schritte lang hielt, stand damit genauso da wie einer, der es über
die gesamte Episode behielt. Der Schaden eines kompromittierten Systems hängt
jedoch von beidem ab: von seinem Wert und von der Verweildauer des Angreifers.
Deshalb erhält der Angreifer je Schritt einen Anteil des Werts aller aktuell
gehaltenen Knoten. Ein fester Betrag je Knoten würde
Arbeitsplatz\-rechner und Produktionsdatenbank gleichsetzen; der wertanteilige Bonus tut das
nicht, und der Einstiegsknoten mit dem Wert null fällt automatisch heraus.

Über die Grundkopplung wirkt dieser Haltebonus unmittelbar auf den
Verteidiger, der damit nicht nur für die Rückeroberung belohnt wird, sondern
auch dafür, sie früh zu erreichen. Genau diese Kopplung bestimmt auch die
Größenordnung des in \Cref{sec:abwehrbonus} eingeführten Abwehrbonus. Der
Verteidiger steht bei einem kompromittierten Knoten stets vor der Wahl,
diesen neu aufzusetzen oder ihn liegen zu lassen und stattdessen einen
bedrohten sauberen Knoten abzuschirmen. \emph{Sauber} heißt dabei, dass der
Angreifer den Knoten nicht besitzt; \emph{bedroht} heißt, dass mindestens ein
Knoten, der ihn laut Topologie erreichen darf, bereits kompromittiert ist, der
Angreifer ihn also unmittelbar angehen kann.

Zwischen diesen beiden Wegen muss das Neuaufsetzen stets die lohnendere
Maßnahme sein,
denn der Verteidiger soll den Angreifer zurückdrängen und nicht neben ihm her
wirtschaften. Andernfalls entstünde ein stilles Einvernehmen: Der Verteidiger
ließe den Angreifer sitzen und unterhielte daneben Sperren, für die er
fortlaufend bezahlt wird, während der Angreifer ungestört seinen Haltebonus
bezieht. Je Schritt kostet ihn das Liegenlassen eines kompromittierten Knotens
$Q$ über die Kopplung $0{,}005 \cdot \text{Wert}(Q)$, und das Abschirmen eines
bedrohten Knotens $Z$ bringt ihm $0{,}00025 \cdot \text{Wert}(Z)$. Damit das
Neuaufsetzen auch im ungünstigsten Fall überwiegt, beim billigsten
kompromittierten Knoten (Wert~10) gegenüber dem wertvollsten bedrohten Ziel
(Wert~100), muss die Haltebonusrate mehr als das Zehnfache der Abwehrbonusrate
betragen. Mit $0{,}005$ gegen $0{,}00025$ liegt das Verhältnis bei 20 und damit
mit Sicherheitsabstand darüber.

\subsection{Abwehrbonus: ein Zustandsreward für Prävention}
\label{sec:abwehrbonus}

Der zentrale Unterschied zum in \Cref{sec:marlon-original} beschriebenen
Original ist ein zustandsbasierter Bonus für erfolgreiche Prävention. Den
Anlass liefert der Referenzlauf: Dort blieb ein verhinderter Zugriff für beide
Seiten folgenlos, weil er dem Angreifer nichts nahm und dem Verteidiger über
die Kopplung deshalb nichts einbrachte (\Cref{sec:referenzlauf}). Ohne ein
Signal, das Prävention sichtbar macht, lässt sich nicht untersuchen, ob eine
Topologie sie begünstigt. Ein
ereignisbasierter Bonus, der jeden abgewehrten Verbindungsversuch einzeln
belohnt, wäre durch Wiederholung ausbeutbar: Der eingefrorene Angreifer läuft in
jeder Episode erneut gegen dieselbe Sperre, und der Verteidiger erhielte bei
jedem Versuch erneut eine Belohnung. Belohnt wird deshalb stattdessen der \emph{Zustand}
je Zeitschritt.

Betrachtet wird jeder gesperrte, saubere Knoten $Z$ mit positivem Wert. Sein
Beitrag bemisst sich am Anteil der Bedrohung an allen Knoten, die $Z$ laut
Topologie erreichen dürfen:
%
\begin{equation}
  f = \frac{\text{kompromittierte Vorgänger von } Z}
           {\text{alle Vorgänger von } Z},
  \qquad
  \text{Beitrag}(Z) = 0{,}00025 \cdot \text{Wert}(Z) \cdot (2f - 1).
  \label{eq:abwehrbonus}
\end{equation}
%
Ist jeder Vorgänger von $Z$ kompromittiert ($f=1$), erhält der Verteidiger den
vollen positiven Beitrag; ist keiner kompromittiert ($f=0$), den vollen
negativen. Dazwischen wird linear interpoliert. Die zweite Hälfte der Formel
ist der eigentliche Kern: Sie verhindert, dass der Verteidiger einfach alles
sperrt und dafür bezahlt wird, denn eine Sperre ohne bestehende Bedrohung
kostet ihn ebenso viel, wie eine Sperre gegen eine vollständige Bedrohung ihm
einbringt.

Eine frühere, hier verworfene Fassung entschied binär, ob \emph{irgendein}
kompromittierter Knoten den Zugangsport von $Z$ überhaupt hätte nutzen dürfen,
unabhängig davon, wie viele saubere Knoten dieselbe Sperre zusätzlich vom
Zugriff ausschloss. In einer Topologie mit hohem Eingangsgrad, etwa
\emph{flat} mit neun eingehenden Kanten je Knoten, zahlte diese Fassung
praktisch jede beliebige Sperre positiv aus, selbst wenn sie überwiegend
legitime Zugänge und nur nebenbei eine tatsächliche Bedrohung traf. Die
Der Anteil $f$ aus \Cref{eq:abwehrbonus} vermeidet das, weil er die Zahl der
Vorgänger im Nenner führt: Je mehr Wege auf einen Knoten führen, desto größer
muss der kompromittierte Teil von ihnen sein, damit sich eine Sperre auszahlt.
Grundlage der Vorgängerbeziehung ist dabei die in \Cref{sec:topomodell}
definierte Adjazenz.

Der Abwehrbonus bleibt gegenüber der \ac{SLA}-Strafe und der Grundkopplung
bewusst klein und soll das auch sein. Der eigentliche Anreiz zur Prävention
entsteht über die Grundkopplung selbst: Ein verhinderter Zugriff ist ein
Knotenwert samt zugehörigem Haltebonus, den der Angreifer nie erhält, und
dieser Effekt gilt unabhängig davon, wie präzise eine einzelne Sperre im Sinne
von \Cref{eq:abwehrbonus} war.

Die Grundkopplung macht das Spiel bis auf zwei Ausnahmen zu einem
Nullsummenspiel: Was der Angreifer gewinnt, verliert der Verteidiger und
umgekehrt. Daraus folgt unmittelbar, dass ein Reimage dem Verteidiger den Wert
des zurückgewonnenen Knotens einbringt, denn der Angreifer verliert ihn, und
die Kopplung spiegelt das. Die beiden Ausnahmen sind der Abwehrbonus selbst
und die \ac{SLA}-Strafe.

\subsection{Verfügbarkeit und \acs{SLA}-Schwelle}
\label{sec:sla}

Die Verfügbarkeit ist die einzige Größe, die dem Verteidiger Grenzen setzt.
Ohne sie könnte er beliebig viele Knoten gleichzeitig neu aufsetzen und
sperren, ohne dafür je zu zahlen. Ihre Berechnung ist deshalb im Einzelnen
anzugeben. \ac{CBS} führt eine Netzverfügbarkeit zwischen null und
eins. Sie ist das mit den Knotengewichten gebildete Mittel über alle Knoten;
im vorliegenden Aufbau tragen alle zehn Knoten das Gewicht~1. Der Beitrag
eines einzelnen Knotens beträgt
%
\begin{equation}
  a_i =
  \begin{cases}
    \dfrac{1 + \sum_{s \in S_i} w_s \cdot [\,s \text{ erreichbar}\,]}
          {1 + \sum_{s \in S_i} w_s} & \text{Knoten in Betrieb,} \\[2ex]
    0 & \text{Knoten wird neu aufgesetzt,}
  \end{cases}
  \label{eq:verfuegbarkeit}
\end{equation}
%
wobei $S_i$ die Dienste des Knotens bezeichnet, $w_s$ deren Gewicht, das hier
durchgängig~1 beträgt, und ein Dienst als erreichbar gilt, wenn er läuft
\emph{und} sein Port nicht durch eine ausdrückliche Sperr-Regel des
Verteidigers blockiert ist. Diese zweite Bedingung war im in
\Cref{sec:marlon-original} beschriebenen Original nicht vorgesehen; ohne sie
wäre Sperren in Bezug auf die Verfügbarkeit vollständig kostenlos gewesen. Ein
Knoten mit einem einzigen Dienst trägt somit $1{,}0$ bei, solange der Dienst
läuft und sein Port offen ist, und $0{,}5$, wenn der Verteidiger ihn selbst
gesperrt hat. Das ist zugleich der Regelfall: Jeder Knoten des Inventars trägt
genau einen Dienst, allein der Einstiegsknoten trägt zwei. Er kommt als Ziel
einer Sperre jedoch ohnehin nicht in Betracht. Eine Sperre trifft nur
eingehende Verbindungen, und der Angreifer sitzt von Beginn an auf ihm und
nutzt ihn als Ausgangspunkt; einen Abwehrbonus trägt eine Sperre auf ihm
ebenfalls nicht ein, denn der Bonus wird für den gesperrten Knoten selbst
berechnet, und dieser ist hier dauerhaft kompromittiert und trägt zudem den
Wert~0.

Entscheidend ist der zweite Fall in \Cref{eq:verfuegbarkeit}: Ein neu
aufgesetzter Knoten trägt \num{15}~Zeitschritte lang nichts bei. Die Schwelle
liegt bei $0{,}60$; sie wird unterschritten, sobald weniger als 60~Prozent der
gewichteten Verfügbarkeit erhalten sind. Bei zehn gleich gewichteten Knoten
bedeutet das: Vier gleichzeitig neu aufgesetzte Knoten ergeben genau $0{,}60$
und lösen noch nicht aus, fünf ergeben $0{,}50$ und damit einen Bruch. Eine
einzelne Sperre verschiebt diese Grenze um $0{,}05$.

Der Verteidiger gerät also nicht durch vorsichtiges Handeln in einen Bruch,
sondern entweder durch ausgedehntes gleichzeitiges Neuaufsetzen oder durch
Sperren in einem Umfang, der einen erheblichen Teil des Netzes gleichzeitig
betrifft. Anders als beim Original ersetzt ein \ac{SLA}-Bruch dabei nicht mehr
den gesamten Schrittwert, sondern ergänzt ihn um $-50$. Solange Brüche selten
blieben, machte diese Unterscheidung wenig Unterschied; seit gesperrte Ports
auf die Verfügbarkeit einwirken, sind Brüche jedoch der Normalfall geworden,
und eine ersetzende Strafe hätte einzelne Episoden vollständig dominiert und
mit allen übrigen unvergleichbar gemacht.

Abschließend sind zwei Randbedingungen des Aufbaus festzuhalten; beide bestimmen, wie
eine Episode verläuft und endet. Eine Episode umfasst höchstens \num{2000} Zeitschritte, wobei der Angreifer
zuerst agiert. Vorzeitig endet sie ausschließlich dann, wenn der Angreifer
sämtliche zehn Knoten übernommen hat. Insbesondere beendet ein
\ac{SLA}-Bruch die Episode nicht, obwohl \ac{CBS} das vorsähe; er wird lediglich
gezählt und bestraft. Damit beruhen alle Matchups auf derselben Datenbasis.

Der Einstiegsknoten gehört dem Angreifer von Beginn an und kann nicht neu
aufgesetzt werden. Die formale Siegbedingung des Verteidigers, den Angreifer
vollständig zu verdrängen, ist damit strukturell unerreichbar. Das ist
beabsichtigt: Wäre der Einstiegsknoten neu aufsetzbar, könnte der Verteidiger
die Episode im ersten Schritt beenden und dafür den Siegbonus von \num{5000}
einstreichen. Zu lernen bliebe allein diese eine Aktion, und alles Übrige,
Sperren wie Abwägen, wäre gegenstandslos. Gemessen wird folglich Eindämmung,
nicht Auslöschung. Für den Angreifer bleibt der verbleibende Fußabdruck
wertlos, denn der Einstiegsknoten trägt den Wert~0.

\section{Konvergenzkriterium}
\label{sec:konvergenz}

Ziel ist festzustellen, ab welcher Episode ein Verteidiger nichts Neues mehr
dazulernt. Diese Episode ist eine der drei Zielgrößen dieser Arbeit
(\Cref{sec:metriken}): Sie misst, wie schnell ein Verteidiger ein stabiles
Verhalten erreicht, und ist damit die Größe, an der sich ein Einfluss der
Topologie auf die Lerneffizienz zeigen müsste. Das Kriterium dient also der
Messung, nicht der Steuerung des Trainings. Die Schwierigkeit liegt darin, ein
solches Plateau von einer zufällig ruhigen Phase zu unterscheiden.

Zwei Bedingungen müssen dafür gleichzeitig erfüllt sein, jeweils über ein
gleitendes Fenster von $W = 15$ Episoden. Die erste betrifft die
\emph{Streuung}: Die Standardabweichung der letzten $W$ Episoden-Rewards,
geteilt durch die bisherige Reward-Spanne des Laufs, darf 0{,}05 nicht
überschreiten. Die zweite betrifft den \emph{Trend}: Der Abstand zwischen dem
Mittelwert der letzten $W$ Episoden und dem der $W$ davor, ebenfalls auf die
Spanne bezogen, darf denselben Wert nicht überschreiten. Formal gilt mit
$r_{t}$ als Reward der Episode~$t$ und $\Delta$ als Spanne aller bisherigen
Rewards:
%
\begin{equation}
  \frac{\sigma\bigl(r_{t-W+1},\dots,r_{t}\bigr)}{\Delta} \le 0{,}05
  \qquad\text{und}\qquad
  \frac{\bigl|\overline{r}_{t-W+1:t} - \overline{r}_{t-2W+1:t-W}\bigr|}{\Delta}
  \le 0{,}05 .
  \label{eq:konvergenz}
\end{equation}
%
Beide Bedingungen müssen über zehn aufeinanderfolgende Episoden gehalten
werden, damit eine einzelne ruhige Phase nicht bereits als Konvergenz gilt.
Als Konvergenzepisode gilt diejenige, in der diese Strecke vollständig ist.
Da zwei Fenster als Historie benötigt werden, können die Bedingungen frühestens
in Episode~30 erstmals zutreffen, und die Strecke ist frühestens in Episode~39
erfüllt. Vor Episode~39 kann ein Lauf demnach nicht als konvergiert gelten.

Die Normierung auf die Spanne statt auf den Mittelwert ist bewusst gewählt. Der
übliche Variationskoeffizient teilt die Standardabweichung durch den
Mittelwert; da der Verteidiger-Reward jedoch um null pendeln kann, würde dieser
Quotient dort unbeschränkt wachsen. Die Spanne bleibt demgegenüber stets ein
belastbarer Bezugswert und erlaubt zudem den Vergleich zwischen Topologien mit
sehr unterschiedlichem Reward-Niveau. Für dieses Kriterium existiert keine
Literaturquelle; es ist empirisch an eigenen Läufen kalibriert.

Anders als noch im Referenzlauf wird dieses Kriterium nicht mehr
verwendet, um ein Training vorzeitig abzubrechen. Alle 120 Läufe erhalten
dasselbe Budget von \num{500000} Schritten; erst nachträglich wird das
Kriterium auf die vollständige Lernkurve angewendet, um die Episode zu
bestimmen, ab der ein Lauf als konvergiert gilt. Dieser Verzicht auf einen laufzeitgesteuerten
Abbruch hat drei Gründe. Der erste ist die Vorgeschichte: In früheren
Läufen, deren Reward-Struktur noch eine andere war, griff das Kriterium
sehr früh, sodass die Läufe nach wenigen Dutzend Episoden endeten und über
den weiteren Verlauf nichts mehr auszusagen war. Die übrigen beiden Gründe
betreffen die Vergleichbarkeit. Nur bei gleichem
Budget ist gewährleistet, dass alle 24 Matchups tatsächlich vergleichbar sind: Ein
früh abgebrochener Lauf hätte strukturell weniger Gelegenheit zum Lernen
gehabt als ein durchgelaufener. Und die Konvergenzepisode ist selbst die
abhängige Variable dieser Arbeit. Würde das Training beim Erreichen des
Kriteriums beendet, ließe sich nicht mehr prüfen, ob die damit unterstellte
Konvergenz tatsächlich gehalten hätte. Die gemessene Konvergenzepisode wäre
dann nichts anderes als der Zeitpunkt, zu dem der Abbruch ausgelöst wurde, und
keine aus dem weiteren Verlauf bestätigte Eigenschaft des Lernens.

Das Kriterium setzt ein unveränderliches Ziel voraus. Diese Annahme ist in
Phase~2 erfüllt, weil der Angreifer eingefroren ist. Für Aufbauten, in denen
beide Agenten gleichzeitig lernen, wäre es nicht zulässig.

\section{Erhobene Metriken}
\label{sec:metriken}

Erhoben wird je Episode eine Reihe von Größen. Sie haben nicht denselben Rang:
Drei von ihnen sind Zielgrößen, an denen die Fragestellung beantwortet wird;
die übrigen dienen dazu, einen beobachteten Unterschied auf sein Zustandekommen
zurückzuführen.

\subsubsection{Zielgrößen}

Der \emph{Konvergenzzeitpunkt} aus \Cref{sec:konvergenz} misst die
Lerneffizienz, also wie schnell ein Verteidiger ein stabiles Verhalten
erreicht, unabhängig davon, auf welchem Niveau. Eine ergänzende Metrik, die
stattdessen einen Leistungsmeilenstein abbildet, nämlich die Episode, ab der
der geglättete Verteidiger-Reward nicht mehr negativ ist, entfällt bewusst:
Unter der Grundkopplung ist dieser Reward strukturell negativ
(\Cref{sec:rewards}).

Die \emph{Crown-Jewel-Quote} gibt den Anteil der Episoden an, in denen der
Angreifer den Datenbankserver übernimmt. Sie misst die Wirksamkeit der
Verteidigung unmittelbar und ist unabhängig von der Reward-Skala, kennt aber
nur zwei Ausgänge.

Feiner auflösend ist deshalb die \emph{Zahl der gehaltenen Knoten}, erhoben als
größte Zahl gleichzeitig kompromittierter Systeme im Verlauf einer Episode. Sie
unterscheidet einen Angreifer, der im Einstiegsknoten steckenbleibt, von einem,
der bis kurz vor das Crown Jewel vordringt, und bildet damit auch Teilerfolge
der Verteidigung ab, die die Crown-Jewel-Quote allein nicht sichtbar macht.

\subsubsection{Größen zur Deutung}

Der \emph{Episoden-Reward beider Agenten} ergibt, über die Episoden
aufgetragen, das Paar von Lernkurven, aus dem der Konvergenzzeitpunkt
abgeleitet wird. Als Wirksamkeitsmaß taugt er nicht: Er gibt an, wie wenig der
Angreifer gewinnt, nicht wie zuverlässig ein Angriff verhindert wurde.
Aufschlussreich ist vor allem der Abstand zwischen beiden Kurven. Der Reward
des Angreifers entsteht allein aus dem Geschehen im Netz, der des Verteidigers
zusätzlich aus Abwehrbonus und \ac{SLA}-Strafe. Weichen die Kurven
voneinander ab, so misst der Abstand demnach, was den Verteidiger sein eigenes
Vorgehen kostet.

Die \emph{Zahl der Schritte mit \ac{SLA}-Bruch} beziffert diesen Anteil
unmittelbar. Sie ist die Gegenrechnung zu jeder Verteidigungsleistung: Eine
niedrige Crown-Jewel-Quote ist wenig wert, wenn sie mit einem über weite
Strecken abgeriegelten Netz erkauft wurde.

Die \emph{Aktionsstatistik} protokolliert, welche Aktionstypen beide Agenten
gewählt haben, getrennt nach gültigen und ungültigen Aktionen. Aus ihr lassen
sich Verhaltensänderungen ablesen, die der Reward allein nicht sichtbar macht,
etwa eine Verschiebung zwischen Sperren und Neuaufsetzen im Verlauf eines
Laufs. Mitgeschrieben werden außerdem der Abwehrbonus je Episode und die Zahl
der dabei abgeschirmten Knoten. Für die Höhe des Rewards ist der Bonus
bedeutungslos, für die Deutung nicht: Er ist die einzige Größe, die Prävention
unmittelbar ausweist, und damit die Probe darauf, ob die in
\Cref{sec:abwehrbonus} beabsichtigte Wirkung überhaupt eintritt.

Alle genannten Größen werden je Episode mitgeschrieben und beschreiben damit
zunächst den Trainingsverlauf. Im folgenden Abschnitt werden dieselben Größen
ein zweites Mal erhoben, dann an einer fertig trainierten Policy.

\section{Wirksamkeitsevaluation}
\label{sec:evaluation}

Das Training beantwortet, wie schnell ein Verteidiger stabil wird und welches
Vorgehen er sich dabei aneignet. Es beantwortet nicht, wie gut der fertig
trainierte Verteidiger ist. Ein Agent
kann auf einem schlechten Niveau ebenso zuverlässig konvergieren wie auf einem
guten. Für die Frage nach der Wirksamkeit ist deshalb ein zweiter Durchgang
nötig, der von der Lerngeschwindigkeit vollständig getrennt ist.

\subsection{Aufbau}

In der Evaluation sind \emph{beide} Agenten eingefroren. Es findet keine
Gewichtsanpassung mehr statt; die fertigen Policies werden nur noch
ausgeführt. Gemessen wird damit das Ergebnis des Trainings, nicht sein
Zustandekommen.

Jeder der 120 Trainingsläufe wird in drei Stufen wiederholt, die sich
ausschließlich im Verteidiger unterscheiden:

\begin{description}
  \item[Kein Verteidiger] Der Angreifer agiert ungehindert. Diese Stufe liefert
    die Obergrenze dessen, was in der jeweiligen Topologie überhaupt zu
    erreichen ist, und dient als Bezugsgröße.
  \item[Zufallsverteidiger] Ein Agent, der in jedem Schritt gleichverteilt aus
    dem Aktionsraum aus \Cref{tab:aktionsraum-def} zieht. Er handelt also,
    ohne etwas gelernt zu haben.
  \item[Trainierter Verteidiger] Das \ac{PPO}-Modell des jeweiligen Laufs.
\end{description}

Der Angreifer ist in allen drei Stufen derselbe und wird aus demselben
Lauf-Ordner geladen wie das Verteidigermodell. Damit unterscheiden sich die
Stufen tatsächlich nur im Verteidiger.

Der Zufallsverteidiger ist dabei die eigentliche Messlatte, nicht das
ungeschützte Netz. Dass ein Verteidiger besser abschneidet als gar keiner,
sagt wenig; die Frage ist, ob das Gelernte über zufälliges Handeln hinausgeht.

Je Zelle werden \num{25} Episoden mit derselben Episodenlänge von höchstens
\num{2000} Schritten gespielt. Bei vier Topologien, sechs Angreifern, fünf
Seeds und drei Stufen ergeben sich \num{360} Zellen und damit \num{9000}
Episoden.

\subsection{Auswertungsgrößen}

Als Hauptgröße dient der \emph{Restanteil}: der Angreifer-Reward mit
Verteidiger, geteilt durch den Angreifer-Reward ohne Verteidiger, jeweils im
selben Matchup. Ein Wert von 100~Prozent bedeutet, dass der Verteidiger nichts
bewirkt, ein Wert nahe null, dass der Angriff weitgehend unterbunden wird.
Gebildet wird er in drei Stufen: zunächst der Median über die Episoden einer
Zelle, dann der Median über die fünf Seeds, und erst danach das Verhältnis.
Berechnet wird er für den trainierten und für den Zufallsverteidiger
gleichermaßen, beide gegen dieselbe Stufe ohne Verteidiger. Die Differenz
beider Restanteile ist der eigentliche Befund, denn sie beziffert, was das
Training über zufälliges Handeln hinaus einbringt.

Auf ein Mittel über die sechs Angreifer wird bewusst verzichtet. Deren
erreichbare Niveaus liegen weit auseinander, und beide naheliegenden
Aggregationen verzerren in unterschiedliche Richtungen: Ein Verhältnis der
Summen wird vom stärksten Angreifer bestimmt, ein Mittel der Einzelverhältnisse
vom schwächsten. Die Ergebnisse werden deshalb je Matchup ausgewiesen.

Ergänzend dienen dieselben beiden Zielgrößen wie im Training, die Zahl der vom
Angreifer gehaltenen Knoten und die Crown-Jewel-Quote. Der Restanteil sagt, um
welchen Faktor der Ertrag des Angreifers schrumpft, aber nicht, woran das
liegt; ein halbierter Reward kann bedeuten, dass der Angreifer auf halbem Weg
stehen blieb, oder dass er zwar durchkam, aber deutlich länger brauchte. Erst
die beiden Zielgrößen entscheiden zwischen diesen Fällen.
